% !TeX root = ../../Book5.tex
% 纲要标题：## 第6章　平坦性、纤维与泛自由性
% 纲要位置：几何学计划纲要.md，第 5220 行。
% 正文状态：初稿；尚未正式审核。

% 纲要内容开始
% 
% > 本章目标与证明范围：把平坦性作为相对几何的核心工具独立建立，完整证明实际使用的判据和泛自由性。纤维维数、矩阵秩、局部自由及平坦性之间的假设逐项落实。
% 
% 纲要内容结束

\chapter{平坦性、纤维与泛自由性}
\label{g5:ch:06}
\BookChapterNumber{g5:ch:06}

\begin{chapterabstract}
	本章研究随基环变化的模与方程族。有限关系判据和有限生成理想判据把平坦性的定义转为可验证的代数条件，局部化则将这些条件送到各个点。Fitting 理想从呈示矩阵的子式中提取与选择无关的信息，区分纤维维数、局部自由性和幂零方向上的关系。泛自由性通过分式域上的标准单项式及有限次清除分母，证明 Noether 整基环上的有限型对象在某个稠密开集上自由，秩可以无限。忠实平坦性进一步使系数扩张能够反映正合性，并为双重基变换中的相容条件提供代数依据。
\end{chapterabstract}

\begin{learninggoals}
	\begin{itemize}
		\item 用有限关系和理想张量判据证明平坦性，并区分整环上的无挠条件与平坦条件。
		\item 从有限呈示计算 Fitting 理想，证明其选择无关性及基变换公式，判定纤维维数和局部自由秩。
		\item 完整执行泛自由性证明中的标准单项式选择和清除分母步骤，说明为何无限秩仍只需一个分母。
		\item 用极大理想和谱映射检验忠实平坦性，在基变换后检测消失、同构与正合性。
		\item 构造平坦曲线族，比较纤维的分支、奇点、长度及维数，并通过双重基变换识别来自原模的元素。
	\end{itemize}
\end{learninggoals}

考虑两组随参数 \(t\) 变化的方程：
\[
	xy=t,\qquad tz=0.
\]
第一组在 \(t\ne0\) 时给出双曲线，在 \(t=0\) 时给出两条相交直线；第二组在 \(t\ne0\) 时迫使 \(z=0\)，在 \(t=0\) 时却失去对 \(z\) 的限制。两者都发生形状变化，但函数之间的关系发生变化的方式不同。在第二组的坐标环里，原来非零的参数 \(t\) 杀死了非零函数 \(z\)；第一组则能选出一组对所有参数都有效的 \(k[t]\) 模基。

\ProseParagraphBreak

第二章已经用张量积保持单射定义平坦性，并证明有限呈示平坦模局部自由。本章从这里继续：怎样只检查有限的关系来验证平坦性，怎样从矩阵中读出局部自由条件，以及在一个任意给定的有限型族中，能否删去一块闭集后统一选择模基。最后一个问题由泛自由性回答；其证明的关键在于先控制关系的首项，而非给分式域上的每一个基向量分别选择分母。

\ProseParagraphBreak

所需基础为第一章的单项式序与多项式消元、第二章的张量积、局部化和有限呈示模，以及第三章的素谱纤维与维数。平坦性的代数工具可对照 \cite[第7节，第45—53页]{Matsumura1989}；泛自由性的一个更一般版本见同书第24节定理24.1，第185—186页。本章采用标准单项式证明，并逐步说明每次有限性假设的用途。

% !TeX root = ../../Book5.tex
% 纲要标题：### 6.1 平坦性的代数判据
% 纲要位置：几何学计划纲要.md，第 5224 行。
% 正文状态：初稿；尚未正式审核。

% 纲要内容开始
% 
% * 张量右正合与左端障碍
% * 有限关系判据
% * 有限生成理想张量判据
% * 局部化及局部检测
% * 无挠和平坦的不同范围
% 
% 纲要内容结束

\section{平坦性的代数判据}
\label{g5:sec:0601}

把一个线性方程组的系数送到另一个环后，原来的解之间可能出现新的关系。张量积精确地实施这一系数变换，而平坦性要求其中原有的单射仍能保留下来。第二章已经给出平坦模的定义及其与有限投射模的联系；现在需要把“检验所有单射”变成可以实际使用的代数条件。关键在于：一个张量元素及其所涉及的关系总是有限的，因此可以先研究有限个系数构成的方程，再把结论用于任意模。

\subsection{张量右正合与左端障碍}
\label{g5:subsec:0601:01}

以下环均为交换含幺环，模均为含幺模。由\cref{g5:prop:tensor-right-exact}，任意短正合列
\[
	0\longrightarrow U\longrightarrow V\longrightarrow W\longrightarrow0
\]
与一个 \(A\) 模 \(M\) 张量后，总得到右正合列
\[
	U\otimes_A M\longrightarrow V\otimes_A M
	\longrightarrow W\otimes_A M\longrightarrow0.
\]
按照\cref{g5:def:flat-module}，\(M\) 平坦恰好表示第一张映射也总为单射。特别地，平坦性保持任意正合列：把相邻同态经各自的像分解，便可逐段应用短正合列的保持性质。

\ProseParagraphBreak

这个左端条件测量的是具体关系。例如\cref{g5:exm:tensor-loses-injection}中，\(A=k[t]\) 上乘以 \(t\) 的单射，张量 \(A/(t)\) 后变成了零映射。作为系数作用的 \(t\) 消失了，原来的嵌入也随之消失。与此相对，自由模、有限投射模以及局部化 \(S^{-1}A\) 都平坦。对任意指标集 \(J\)，公式
\[
	N\otimes_A\Bigl(\bigoplus_{j\in J}A\Bigr)
	\cong\bigoplus_{j\in J}N
\]
把张量的有限和逐坐标展开，说明任意秩的自由模都保持单射；投射直和分量也保持单射。局部化的结论则来自\cref{g5:thm:localization-exact,g5:prop:localization-tensor}。

\begin{definition}\label{g5:def:flat-homomorphism}
	设 \(\varphi:A\to B\) 为交换含幺环同态。若 \(B\) 经 \(\varphi\) 成为平坦 \(A\) 模，就称 \(\varphi\) 为\term[pingtantongtai]{平坦同态}[flat homomorphism]，也称 \(B\) 为平坦 \(A\) 代数。
\end{definition}

平坦同态不必单射。例如环的投影 \(k\times k\to k\) 是平坦同态，因为目标是源环作为模的一个直和分量。这与平坦模定义中的“保持单射”不矛盾：定义要求保持的是与该模张量的任意模同态，并没有要求结构同态 \(A\to B\) 自身单射。

\ProseParagraphBreak

为了处理所有模中的包含关系，可以先把较大的模写成自由模的商。下面的归约说明，即使检验对象限制为有限秩自由模中的子模，仍足以检测平坦性。子模本身无需有限生成。

\begin{lemma}\label{g5:lem:flat-submodule-free-test}
	设 \(A\) 为交换含幺环，\(M\) 为 \(A\) 模。则 \(M\) 平坦，当且仅当对每个有限整数 \(r\geq0\) 及每个子模 \(H\subseteq A^r\)，自然映射
	\[
		H\otimes_A M\longrightarrow A^r\otimes_A M
	\]
	都是单射。
\end{lemma}
\begin{proof}
	必要性就是平坦性的定义。为证充分性，先把假设扩展到任意自由模 \(F=\bigoplus_{j\in J}Ae_j\) 的子模 \(H\)。若
	\[
		z=\sum_{i=1}^n h_i\otimes m_i\in H\otimes_A M
	\]
	在 \(F\otimes_A M\) 中为零，所有 \(h_i\) 只涉及有限多个基向量。取这些基向量张成的有限自由直和分量 \(F_0\subseteq F\)，并令 \(H_0=H\cap F_0\)。同一个表达式给出 \(z_0\in H_0\otimes_A M\)。因 \(F_0\) 是直和分量，\(F_0\otimes_A M\to F\otimes_A M\) 仍单射，故 \(z_0\) 在 \(F_0\otimes_A M\) 中为零。假设使 \(z_0=0\)，从而 \(z=0\)。这证明了任意自由模中的子模检验。

	现在取任意包含 \(U\subseteq V\)。选自由模满射 \(\pi:F\to V\)，令
	\[
		K=\ker\pi,\qquad H=\pi^{-1}(U).
	\]
	则 \(H/K\cong U\)、\(F/K\cong V\)。张量右正合性给出交换图
	\[
		\begin{tikzcd}[column sep=small]
			K\otimes_A M\arrow[r]\arrow[d,equal]&
			H\otimes_A M\arrow[r]\arrow[d]&
			U\otimes_A M\arrow[r]\arrow[d]&0\\
			K\otimes_A M\arrow[r]&
			F\otimes_A M\arrow[r]&
			V\otimes_A M\arrow[r]&0 .
		\end{tikzcd}
	\]
	两行正合，中间竖箭头由刚证的自由模情形为单射。

	若 \(z\in U\otimes_A M\) 在下行中变成零，先把它提升为 \(w\in H\otimes_A M\)。\(w\) 在 \(F\otimes_A M\) 中的像属于 \(K\otimes_A M\) 的像，故可选 \(q\in K\otimes_A M\) 给出这个像。从 \(w\) 减去 \(q\) 在上行中的像后，所得元素在 \(F\otimes_A M\) 中为零，由中间竖箭头单射可知它为零。因此 \(w\) 本来就在上行左箭头的像中，进而 \(z=0\)。于是任意 \(U\subseteq V\) 张量后仍单射，\(M\) 平坦。
\end{proof}

这个证明把任意单射中的一个待检验元素，逐步放到了有限个自由坐标中。下一步便是把“在自由模中映为零”写成有限线性方程组。

\subsection{有限关系判据}
\label{g5:subsec:0601:02}

给定 \(a_1,\ldots,a_n\in A\)，标量解
\[
	(a_1,\ldots,a_n)
	\begin{pmatrix}b_1\\ \vdots\\ b_n\end{pmatrix}=0
\]
可以与任意 \(y\in M\) 相乘，得到模中的解 \((b_1y,\ldots,b_ny)\)。把有限多个这样的解相加，仍得到模中的解。平坦性的有限关系判据断言：平坦模中的每一个有限关系，都能由这种方式获得。

\begin{theorem}\label{g5:thm:flat-equational-criterion}
	设 \(A\) 为交换含幺环，\(M\) 为任意 \(A\) 模。下列条件等价：
	\begin{equivalences}
		\item \(M\) 为平坦 \(A\) 模。
		\item 对任意 \(n\geq1\)、\(a_1,\ldots,a_n\in A\) 及 \(x_1,\ldots,x_n\in M\)，若
		\[
			\sum_{j=1}^n a_jx_j=0,
		\]
		则存在有限个 \(y_1,\ldots,y_s\in M\) 和 \(b_{j\ell}\in A\)，使
		\[
			x_j=\sum_{\ell=1}^s b_{j\ell}y_\ell,\qquad
			\sum_{j=1}^n a_jb_{j\ell}=0
			\quad(1\leq\ell\leq s).
		\]
		\item 对任意有限矩阵 \(C=(a_{ij})\in\operatorname{Mat}_{r\times n}(A)\)，若列向量 \(x=(x_j)\in M^n\) 满足 \(Cx=0\)，则存在有限矩阵 \(E=(b_{j\ell})\in\operatorname{Mat}_{n\times s}(A)\) 和 \(y=(y_\ell)\in M^s\)，使
		\begin{equation}\label{g5:eq:flat-relation-factorization}
			x=Ey,\qquad CE=0.
		\end{equation}
	\end{equivalences}
	这里允许 \(s=0\)，此时相应的和为零。
\end{theorem}
\begin{proof}
	先设 \(M\) 平坦。把矩阵 \(C\) 看作同态 \(A^n\to A^r\)，令 \(K=\ker C\)。平坦性保持
	\[
		K\longrightarrow A^n\xrightarrow{C}A^r
	\]
	在中间的正合性，故 \(x\in\ker(C\otimes1_M)\) 来自某个 \(K\otimes_A M\) 的元素。张量元素可写成有限和
	\[
		\sum_{\ell=1}^s v_\ell\otimes y_\ell,\qquad
		v_\ell=(b_{1\ell},\ldots,b_{n\ell})^{\mathsf T}\in K.
	\]
	以 \(v_\ell\) 为列组成矩阵 \(E\)，就得到 \(x=Ey\) 及 \(CE=0\)。因此第一条件蕴含第三条件；取 \(r=1\)，又得到第二条件。

	接下来证明，第二条件中的一个方程已经足以处理任意有限个方程。对方程数 \(r\) 归纳。\(r=1\) 正是第二条件。若 \(r>1\)，先对第一行方程应用该条件，把 \(x\) 写成 \(x=E_1y\)，并使 \(C\) 的第一行乘 \(E_1\) 为零。把这个表达代入其余 \(r-1\) 行，得到关于 \(y\) 的新方程组
	\[
		C' E_1y=0,
	\]
	其中 \(C'\) 由 \(C\) 的后 \(r-1\) 行组成。归纳假设给出 \(y=E_2z\)、\(C'E_1E_2=0\)。于是 \(E=E_1E_2\) 同时满足 \(x=Ez\) 和 \(CE=0\)。若某一步已经得到 \(x=0\)，取空矩阵即可。因此第二条件蕴含第三条件。

	最后设第三条件成立。按\cref{g5:lem:flat-submodule-free-test}，只需检验任意子模 \(H\subseteq A^r\) 张量后仍嵌入。设
	\[
		z=\sum_{j=1}^n h_j\otimes x_j\in H\otimes_A M
	\]
	在 \(A^r\otimes_A M=M^r\) 中为零。写 \(h_j\) 的第 \(i\) 个坐标为 \(a_{ij}\)，组成矩阵 \(C=(a_{ij})\)。映为零恰好是 \(Cx=0\)。由第三条件选 \(x=Ey\)、\(CE=0\)，便有
	\[
		z=\sum_{\ell=1}^s
		\left(\sum_{j=1}^n b_{j\ell}h_j\right)\otimes y_\ell.
	\]
	每个括号中的元素，其全部坐标由 \(CE=0\) 为零；它在子模 \(H\) 中也为零。因此 \(z=0\)，所需单射成立，\(M\) 平坦。
\end{proof}

有限关系判据中的“有限”约束的是每次被检验的关系，而不是整个模的生成元数。它也没有要求 \(\ker C\) 有限生成：每个具体的张量只取用了其中有限多个向量。这个区别使判据能够用于非 Noether 环和无限生成模。\cref{g5:thm:flat-equational-criterion}对应\cite[第7节，定理7.6，第49--50页]{Matsumura1989}。

\ProseParagraphBreak

以一个系数为例，若 \(ax=0\)，判据要求
\[
	x=\sum_\ell b_\ell y_\ell,\qquad ab_\ell=0.
\]
因此平坦模中被 \(a\) 杀死的元素，只能由环中本来被 \(a\) 杀死的系数生成。多系数关系则同时约束这些元素如何组合，这比逐个检验一个标量的乘法要强。

\subsection{有限生成理想张量判据}
\label{g5:subsec:0601:03}

给定理想 \(I\subseteq A\)，自然映射
\[
	\mu_I:I\otimes_A M\longrightarrow M,\qquad
	a\otimes m\longmapsto am
\]
的像总是 \(IM\)。所以这张映射是否单射，等价于：乘积之间的零关系是否已经由张量的平衡关系解释。若 \(I=(a_1,\ldots,a_n)\)，它恰好集中检验带这些系数的有限关系。

\begin{theorem}\label{g5:thm:flat-ideal-criterion}
	设 \(A\) 为交换含幺环，\(M\) 为任意 \(A\) 模。下列条件等价：
	\begin{equivalences}
		\item \(M\) 为平坦 \(A\) 模；
		\item 对每个有限生成理想 \(I\subseteq A\)，乘法映射
		\[
			\mu_I:I\otimes_A M\longrightarrow M
		\]
		为单射；
		\item 对每个理想 \(I\subseteq A\)，\(\mu_I\) 都为单射。
	\end{equivalences}
	这些条件成立时，\(\mu_I\) 自然识别 \(I\otimes_A M\) 与 \(IM\)。
\end{theorem}
\begin{proof}
	平坦性把包含 \(I\hookrightarrow A\) 变成单射
	\(I\otimes_A M\to A\otimes_A M=M\)，所以第一条件蕴含第三条件，第三条件也蕴含第二条件。

	设第二条件成立。任取有限关系 \(\sum_{j=1}^n a_jx_j=0\)，令
	\[
		I=(a_1,\ldots,a_n),\qquad
		q:A^n\longrightarrow I,\quad e_j\longmapsto a_j,\qquad
		K=\ker q.
	\]
	元素 \(z=\sum_j a_j\otimes x_j\in I\otimes_A M\) 满足 \(\mu_I(z)=0\)，由假设 \(z=0\)。张量右正合性使
	\[
		K\otimes_A M\longrightarrow M^n
		\xrightarrow{q\otimes1} I\otimes_A M\longrightarrow0
	\]
	正合。因此列向量 \((x_j)\) 是某个有限和
	\(\sum_\ell v_\ell\otimes y_\ell\)、\(v_\ell\in K\)，的像。写
	\(v_\ell=(b_{1\ell},\ldots,b_{n\ell})^{\mathsf T}\)，便有
	\[
		x_j=\sum_\ell b_{j\ell}y_\ell,\qquad
		\sum_j a_jb_{j\ell}=0.
	\]
	这正是有限关系判据，故 \(M\) 平坦。最后，\(\mu_I\) 的像本来就是 \(IM\)，其单射性给出所述识别。
\end{proof}

这里足够检验的是有限生成理想，而不是要求环的全部理想都有限生成。这一判据见\cite[第7节，定理7.7，第50--51页]{Matsumura1989}。实际使用时，发现某个 \(I\otimes_A M\to M\) 有非零核就能否定平坦性；反方向则必须处理所有有限生成理想。

\ProseParagraphBreak

当 \(I=(a)\) 时，取一个系数的有限关系判据还给出如下等式：
\begin{equation}\label{g5:eq:flat-annihilator}
	(0:_M a)=\operatorname{Ann}_A(a)\,M
	\qquad\text{若 \(M\) 为平坦 \(A\) 模}.
\end{equation}
左端表示被 \(a\) 杀死的模元素，右端表示由被 \(a\) 杀死的环元素作用于 \(M\) 后生成的子模。右端包含于左端对任何模都成立；另一包含由上一小节对 \(ax=0\) 的分解得到。对于没有零因子的 \(a\)，右端为零，因而乘以 \(a\) 在平坦模上也没有核。

\subsection{局部化及局部检测}
\label{g5:subsec:0601:04}

第二章的\cref{g5:prop:flat-local-detection}已经证明：一个 \(A\) 模 \(M\) 平坦，当且仅当全部 \(M_{\mathfrak p}\) 在 \(A_{\mathfrak p}\) 上平坦，而且只查极大理想便足够。其证明先取任意张量映射的核，再用局部化的正合性和模消失的局部检测。为了研究一个 \(A\) 代数 \(B\) 上的模，还需要允许在 \(\operatorname{Spec}B\) 上逐点检验相对于 \(A\) 的平坦性。先记录改变系数环时所用的基本同构。

\begin{proposition}\label{g5:prop:flat-base-change}
	设 \(A\to B\) 为交换含幺环同态。
	\begin{enumerate}
		\item 若 \(M\) 为平坦 \(A\) 模，则 \(B\otimes_A M\) 为平坦 \(B\) 模。
		\item 若 \(B\) 为平坦 \(A\) 模，且 \(N\) 为平坦 \(B\) 模，则 \(N\) 经标量限制成为平坦 \(A\) 模。
	\end{enumerate}
	因此平坦同态在任意基变换及复合下保持平坦。
\end{proposition}
\begin{proof}
	对任意 \(B\) 模 \(U\)，有自然同构
	\[
		U\otimes_B(B\otimes_A M)\cong U\otimes_A M,\qquad
		u\otimes(b\otimes m)\longmapsto bu\otimes m,
	\]
	逆映射为 \(u\otimes m\mapsto u\otimes(1\otimes m)\)。一个 \(B\) 模单射也为 \(A\) 模单射，故右侧由 \(M\) 的平坦性保持单射，得到第一断言。这里不要求 \(B\) 平坦于 \(A\)。

	对任意 \(A\) 模 \(U\)，又有
	\[
		U\otimes_A N\cong(U\otimes_A B)\otimes_B N,
	\]
	同构把 \((u\otimes b)\otimes n\) 送到 \(u\otimes bn\)。若 \(B\) 平坦于 \(A\)，第一步张量保持单射；若 \(N\) 平坦于 \(B\)，第二步也保持单射，故得到第二断言。将第一断言用于一个平坦代数的系数变换，将第二断言用于两步平坦环同态，便得到最后的结论。
\end{proof}

特别地，\(S^{-1}M\cong S^{-1}A\otimes_A M\) 在 \(M\) 平坦时为平坦 \(S^{-1}A\) 模。它同时也是平坦 \(A\) 模，因为 \(S^{-1}A\) 在 \(A\) 上平坦，再用第二断言即可。基变换及复合性质可对照\cite[第7节，第46页]{Matsumura1989}。

\begin{proposition}\label{g5:prop:relative-flat-localization}
	设 \(A\to B\) 为交换含幺环同态，\(M\) 为 \(B\) 模且作为 \(A\) 模平坦，\(T\subseteq B\) 为含 \(1\) 的乘法集。则 \(T^{-1}M\) 作为 \(A\) 模仍平坦。
\end{proposition}
\begin{proof}
	对任意 \(A\) 模 \(U\)，\(U\otimes_A M\) 通过第二因子成为 \(B\) 模，并有自然同构
	\[
		T^{-1}(U\otimes_A M)\cong U\otimes_A T^{-1}M,\qquad
		\frac{u\otimes m}{t}\longmapsto u\otimes\frac mt.
	\]
	它也可由
	\(T^{-1}B\otimes_B(U\otimes_A M)\) 的张量结合律得到。对任意 \(A\) 模单射 \(U\to V\)，先与 \(M\) 张量得到 \(B\) 模单射，再按 \(T\) 局部化仍为单射。用上述同构识别结果，就证明 \(T^{-1}M\) 平坦于 \(A\)。
\end{proof}

\begin{proposition}\label{g5:prop:flat-local-test}
	设 \(\varphi:A\to B\) 为交换含幺环同态，\(M\) 为任意 \(B\) 模。对 \(\mathfrak q\in\operatorname{Spec}B\)，令
	\(\mathfrak p=\varphi^{-1}(\mathfrak q)\)，并通过 \(A_{\mathfrak p}\to B_{\mathfrak q}\) 把 \(M_{\mathfrak q}\) 看作 \(A_{\mathfrak p}\) 模。下列条件等价：
	\begin{equivalences}
		\item \(M\) 作为 \(A\) 模平坦；
		\item 对每个 \(\mathfrak q\in\operatorname{Spec}B\)，\(M_{\mathfrak q}\) 为平坦 \(A_{\mathfrak p}\) 模；
		\item 对每个极大理想 \(\mathfrak q\subseteq B\)，\(M_{\mathfrak q}\) 为平坦 \(A_{\mathfrak p}\) 模。
	\end{equivalences}
\end{proposition}
\begin{proof}
	若 \(M\) 为平坦 \(A\) 模，由\cref{g5:prop:relative-flat-localization}，\(M_{\mathfrak q}\) 仍平坦于 \(A\)。对任意 \(A_{\mathfrak p}\) 模 \(U\)，有
	\[
		U\otimes_A M_{\mathfrak q}
		\cong U\otimes_{A_{\mathfrak p}}M_{\mathfrak q}.
	\]
	确实，分母 \(s\in A\setminus\mathfrak p\) 已在两因子上可逆，张量中的关系
	\((s^{-1}u)\otimes m=u\otimes(s^{-1}m)\)
	可通过把原有的 \(A\) 平衡关系应用于 \(s^{-1}u,s^{-1}m\) 得到。因此一个 \(A_{\mathfrak p}\) 模单射与 \(M_{\mathfrak q}\) 张量后仍单射，得到第二条件。第二条件当然蕴含第三条件。

	反过来，设第三条件成立。给定任意 \(A\) 模单射 \(U\hookrightarrow V\)，记
	\[
		K=\ker\bigl(U\otimes_A M\longrightarrow V\otimes_A M\bigr).
	\]
	这是一个 \(B\) 模。在任意极大理想 \(\mathfrak q\subseteq B\) 处局部化，局部化正合性及张量同构给出
	\[
		K_{\mathfrak q}
		\cong\ker\left(
		U_{\mathfrak p}\otimes_{A_{\mathfrak p}}M_{\mathfrak q}
		\longrightarrow
		V_{\mathfrak p}\otimes_{A_{\mathfrak p}}M_{\mathfrak q}
		\right).
	\]
	这里先将 \(B\) 的分母移到 \(M\) 一侧，再将 \(A\setminus\mathfrak p\) 的分母移到第一因子，便得到所用同构。原单射局部化为 \(U_{\mathfrak p}\hookrightarrow V_{\mathfrak p}\)，而 \(M_{\mathfrak q}\) 按假设平坦，所以 \(K_{\mathfrak q}=0\)。将\cref{g5:prop:local-detection}用于 \(B\) 模 \(K\)，得 \(K=0\)。于是 \(M\) 平坦于 \(A\)。
\end{proof}

这就是\cite[第7节，定理7.1，第46--47页]{Matsumura1989}的相对局部检测。它检验的是 \(M_{\mathfrak q}\) 在 \(A_{\mathfrak p}\) 上的平坦性，并非在 \(B_{\mathfrak q}\) 上的平坦性；后一个性质属于另一组系数关系。此外，\(B\) 的极大理想收缩到 \(A\) 后未必极大，例如 \(\mathbb Z\to\mathbb Q\) 中的 \((0)\) 收缩仍为 \((0)\)。陈述中的 \(A_{\mathfrak p}\) 因而必须保留实际的收缩素理想。

\ProseParagraphBreak

若主开集 \(D(f_i)\) 覆盖 \(\operatorname{Spec}A\)，则 \(M\) 平坦当且仅当每个 \(M_{f_i}\) 平坦于 \(A_{f_i}\)。必要性来自局部化；充分性是在任意 \(\mathfrak p\) 处选 \(f_i\notin\mathfrak p\)，再局部化得到平坦 \(A_{\mathfrak p}\) 模，最后应用第二章的局部检测。这一叙述无需模有限生成，也不要求给定覆盖为有限覆盖。

\subsection{无挠和平坦的不同范围}
\label{g5:subsec:0601:05}

平坦性首先排除了某些标量产生的新零关系。若基环为整环，每个非零标量在环上都给出单射，所以它在平坦模上也必须给出单射。这就是无挠性。但在多个系数同时出现时，还可能发生一个标量看不出的障碍。

\begin{definition}\label{g5:def:torsion-free-module}
	设 \(A\) 为整环，\(M\) 为 \(A\) 模。若对任意 \(0\ne a\in A\) 和 \(m\in M\)，等式 \(am=0\) 都蕴含 \(m=0\)，则称 \(M\) 为\term[wuniaomo]{无挠模}[torsion-free module]。
\end{definition}

\begin{proposition}\label{g5:prop:flat-torsion-free}
	设 \(A\) 为整环，\(M\) 为平坦 \(A\) 模。则 \(M\) 无挠。
\end{proposition}
\begin{proof}
	对任意 \(0\ne a\in A\)，乘法映射 \(A\xrightarrow{a}A\) 为单射。与 \(M\) 张量后得到 \(M\xrightarrow{a}M\)，由平坦性仍为单射。
\end{proof}

\begin{example}\label{g5:exm:torsion-free-nonflat-ideal}
	设 \(k\) 为任意域，\(A=k[x,y]\)，\(I=(x,y)\subseteq A\)。作为整环的子模，\(I\) 无挠；但 \(I\) 不是平坦 \(A\) 模。

	为证明不平坦，考察
	\[
		\mu_I:I\otimes_A I\longrightarrow I,\qquad
		f\otimes g\longmapsto fg.
	\]
	元素
	\[
		z=x\otimes y-y\otimes x
	\]
	映到 \(xy-yx=0\)。我们还必须确认 \(z\) 在张量积中非零。令 \(\mathfrak m=(x,y)\)，有自然映射
	\[
		I\otimes_A I\longrightarrow
		(I/\mathfrak mI)\otimes_k(I/\mathfrak mI).
	\]
	由于 \(\mathfrak mI=(x,y)^2\)，商模 \(I/\mathfrak mI\) 以 \(\bar x,\bar y\) 为 \(k\) 基。\(z\) 的像为
	\[
		\bar x\otimes\bar y-\bar y\otimes\bar x,
	\]
	它是两个不同的张量基向量之差，在任意特征下都非零。所以 \(z\ne0\)，\(\mu_I\) 不单射，由\cref{g5:thm:flat-ideal-criterion}知 \(I\) 不平坦。这个障碍由两个生成元之间的关系 \(yx-xy=0\) 产生，逐个检验非零标量的单射性看不到它。
\end{example}

若所有理想都由一个元素生成，理想判据就把平坦性完全归约到一个标量的单射性。

\begin{proposition}\label{g5:prop:torsion-free-pid-flat}
	设 \(A\) 为主理想整环，\(M\) 为任意 \(A\) 模。则 \(M\) 平坦当且仅当 \(M\) 无挠。特别地，该结论对任意离散赋值环成立，且不要求 \(M\) 有限生成。
\end{proposition}
\begin{proof}
	平坦蕴含无挠已经证明。反过来，设 \(M\) 无挠。任意理想 \(I\subseteq A\) 都为主理想。若 \(I=0\)，\(I\otimes_A M\to M\) 自然单射；若 \(I=aA\)、\(a\ne0\)，则
	\[
		A\xrightarrow{\ \sim\ }I,\qquad r\longmapsto ar
	\]
	是模同构。在由此得到的识别 \(I\otimes_A M\cong M\) 下，乘法映射 \(\mu_I\) 正是 \(M\xrightarrow{a}M\)。无挠性使它单射，因此理想判据给出平坦性。DVR 为主理想整环，已由\cref{g5:thm:dvr-characterizations}证明，故适用同一结论。
\end{proof}

例如 \(\mathbb Q\) 作为 \(\mathbb Z\) 模无挠，因而平坦；它不是自由 \(\mathbb Z\) 模。确实，乘以 \(2\) 在 \(\mathbb Q\) 上为满射，而在非零自由 \(\mathbb Z\) 模上不满。更一般地，DVR \(R\) 的分式域 \(K\) 平坦于 \(R\)，却不是自由 \(R\) 模：若 \(\pi\) 为统一化元，则 \(K/\pi K=0\)，而任何非零自由 \(R\) 模模去 \(\pi\) 后都有非零剩余向量空间。无限生成的平坦模因而仍有别于自由模。

\begin{corollary}\label{g5:cor:dedekind-torsion-free-flat}
	设 \(A\) 为 Dedekind 整环，允许 \(A\) 为域，\(M\) 为任意 \(A\) 模。则 \(M\) 平坦当且仅当 \(M\) 无挠。
\end{corollary}
\begin{proof}
	只需证明无挠蕴含平坦。无挠性在整环的局部化下保持：若
	\((a/s)(m/t)=0\) 且 \(a/s\ne0\)，则某个被反演的非零元素 \(u\) 满足 \(uam=0\)；\(ua\ne0\)，故 \(m=0\)。因此每个 \(M_{\mathfrak p}\) 都无挠于 \(A_{\mathfrak p}\)。

	由\cref{g5:thm:dedekind-local}，非零素理想处的 \(A_{\mathfrak p}\) 为 DVR，故上一命题说明 \(M_{\mathfrak p}\) 平坦。零素理想处的局部环是分式域，其上所有模都是自由向量空间，也平坦。最后用\cref{g5:prop:flat-local-detection}得到 \(M\) 平坦。若 \(A\) 本来是域，直接使用向量空间的自由性即可。
\end{proof}

Dedekind 整环上的这一结论也见\cite[第11节，习题11.8，第86页]{Matsumura1989}。第四章的局部 DVR 刻画由此把无挠性与平坦性联系起来，而二维多项式环上的理想 \((x,y)\) 已显示这种联系不能延伸到任意整环。一般情况下，主理想只检验单个标量；完整的平坦性需要保留有限生成理想所编码的全部系数关系。
% 纲要条目排为编号小节标题，并进入目录与书签；说明段仍保留于纲要注释中。

% !TeX root = ../../Book5.tex
% 纲要标题：### 6.2 有限呈示与 Fitting 理想
% 纲要位置：几何学计划纲要.md，第 5232 行。
% 正文状态：初稿；尚未正式审核。

% 纲要内容开始
% 
% * 呈示子式和 Fitting 理想
% * 与呈示选择无关及基变换
% * 剩余域上的秩条件
% * 局部自由秩的理想判据
% * 幂零基底上恒定纤维维数的反例
% 
% 纲要内容结束

\section{有限呈示与 Fitting 理想}
\label{g5:sec:0602}

第二章把有限呈示模写成有限生成元模去有限条关系，并用关系矩阵计算纤维维数。矩阵的秩由子式控制，但矩阵本身会随生成元与关系的选择而改变。本节要从这些子式中提取模的内在理想：它们既能给出纤维维数跳变的位置，也能判断关系在一个邻域内是否真正消失。后一项判断尤其重要，因为在剩余域中为零的系数，在原环中仍可能是非零幂零元。

\subsection{呈示子式和 Fitting 理想}
\label{g5:subsec:0602:01}

设 \(A\) 为交换含幺环，\(M\) 为有限呈示 \(A\) 模，选定呈示
\begin{equation}\label{g5:eq:fitting-presentation}
	A^m\xrightarrow{H}A^n\longrightarrow M\longrightarrow0.
\end{equation}
矩阵 \(H\) 有 \(n\) 行、\(m\) 列；第 \(j\) 列是第 \(j\) 条关系在所选 \(n\) 个生成元中的系数。关系之间可以继续有关系，因此左端不要求单射。将这个呈示张量到一个域 \(K\) 后，余核维数为 \(n-\operatorname{rank}H_K\)。于是，当我们关心余核至多需要 \(i\) 个方向时，应检查 \(n-i\) 阶子式，而不是固定阶数的子式。

\begin{definition}\label{g5:def:fitting-ideal}
	设 \(A\) 为交换含幺环，\(H\) 为一个 \(n\times m\) 矩阵。对整数 \(q\geq1\)，记 \(I_q(H)\) 为全部 \(q\times q\) 子式生成的理想；若 \(q>\min(n,m)\)，规定 \(I_q(H)=0\)。此外规定 \(I_0(H)=A\)。

	对有限呈示模 \(M\)、呈示\eqref{g5:eq:fitting-presentation}及整数 \(i\geq0\)，定义第 \(i\) 个\term[Fittinglixiang]{Fitting 理想}[Fitting ideal]为
	\[
		\operatorname{Fitt}_i(M)=
		\begin{cases}
			I_{n-i}(H),&0\leq i<n,\\
			A,&i\geq n.
		\end{cases}
	\]
	下节将证明该理想与所选有限呈示无关。为统一秩零的公式，另约定 \(\operatorname{Fitt}_{-1}(M)=0\)。
\end{definition}

这些理想以 Hans Fitting 命名。
\footnote{Hans Fitting（1906--1938），德国数学家，研究群的分解、幂零群及有限生成模的行列式理想。}下标记录的是商模允许保留的生成方向数，所以子式的阶数随下标增大而减小。沿一行展开行列式，任何 \(q+1\) 阶子式都是若干 \(q\) 阶子式的系数倍之和，故
\[
	I_{q+1}(H)\subseteq I_q(H),\qquad
	0\subseteq\operatorname{Fitt}_0(M)
	\subseteq\operatorname{Fitt}_1(M)\subseteq\cdots\subseteq A.
\]
由于矩阵有限，每个 Fitting 理想有限生成，而且从某个下标开始恒等于 \(A\)。

\begin{example}\label{g5:exm:fitting-free-cyclic}
	设 \(A\) 为交换含幺环。自由模 \(M=A^r\) 可取没有关系的呈示 \(0\to A^r\to M\to0\)，所以
	\[
		\operatorname{Fitt}_i(A^r)=
		\begin{cases}0,&0\leq i<r,\\A,&i\geq r.\end{cases}
	\]
	这也包括 \(r=0\)：零模的全部非负下标 Fitting 理想都是 \(A\)。

	若 \(J=(a_1,\ldots,a_m)\subseteq A\)，循环模 \(A/J\) 的关系矩阵是单行矩阵 \((a_1\ \cdots\ a_m)\)，故
	\[
		\operatorname{Fitt}_0(A/J)=J,
		\qquad\operatorname{Fitt}_i(A/J)=A\quad(i\geq1).
	\]
	因此这些理想保存的是关系的实际系数；即使某个 \(a_j\) 幂零，它仍可能出现在 \(\operatorname{Fitt}_0\) 中。
\end{example}

\subsection{与呈示选择无关及基变换}
\label{g5:subsec:0602:02}

换一组关系只改变关系列的生成方式；换一组模生成元则还会改变矩阵的行数。证明呈示无关性时必须同时处理这两种变化。先记录所需的矩阵事实，它们只使用行列式的多重线性，并不要求底环为域或为 Noether 环。

\begin{lemma}\label{g5:lem:determinantal-matrix-operations}
	设矩阵的系数属于交换含幺环 \(A\)。当乘积有定义时，对 \(q\geq1\) 有
	\[
		I_q(UH)\subseteq I_q(H),\qquad
		I_q(HV)\subseteq I_q(H).
	\]
	若 \(U,V\) 分别为可逆方阵，则 \(I_q(UHV)=I_q(H)\)。另对任意矩阵 \(H\) 及 \(q\geq1\)，有
	\[
		I_q\!\begin{pmatrix}H&0\\0&1\end{pmatrix}
		=I_q(H)+I_{q-1}(H)=I_{q-1}(H).
	\]
\end{lemma}
\begin{proof}
	\(UH\) 的每一行都是 \(H\) 各行的线性组合。对选定的 \(q\) 行和 \(q\) 列，逐行展开行列式，得到 \(H\) 的同阶子式的线性组合：选到重复行的项为零，互异行的项按行次序调整符号后就是相应子式。因此得到第一个包含。逐列展开证明第二个包含。若矩阵可逆，再对 \(U^{-1}\)、\(V^{-1}\) 用同样的包含，便得等号。

	对块矩阵，一个子式若同时选中最后一行和最后一列，就等于 \(H\) 的某个 \(q-1\) 阶子式；两者都不选时，就是 \(H\) 的 \(q\) 阶子式；只选其一时，该子矩阵有一行或一列全零。所有子式因此恰好生成 \(I_q(H)+I_{q-1}(H)\)。由沿一行展开得到的包含 \(I_q(H)\subseteq I_{q-1}(H)\)，结论成立。\(q=1\) 时新增的右下角 \(1\) 生成单位理想，也与 \(I_0(H)=A\) 的约定一致。
\end{proof}

可逆行列变换及单位块消元的这些性质，可对照\cite[第13节附录，第103页]{Matsumura1989}。下面进一步证明，随着自由模的秩一起调整子式阶数，便得到只依赖 \(M\) 的理想。

\begin{theorem}\label{g5:thm:fitting-invariance}
	设 \(A\) 为交换含幺环，\(M\) 为有限呈示 \(A\) 模。\cref{g5:def:fitting-ideal}给出的 \(\operatorname{Fitt}_i(M)\) 与全部有限生成元、关系生成元及有限呈示的选择无关。因此同构的有限呈示模具有相同的 Fitting 理想。
\end{theorem}
\begin{proof}
	先固定满射 \(\pi:A^n\to M\)，只更换关系生成元。若 \(H:A^m\to A^n\) 与 \(H':A^{m'}\to A^n\) 的像都是 \(\ker\pi\)，则 \(H'\) 的每一列都是 \(H\) 各列的线性组合，反之亦然。因此存在矩阵 \(C,D\) 使
	\[
		H'=HC,\qquad H=H'D.
	\]
	由上个引理，\(I_q(H')\subseteq I_q(H)\subseteq I_q(H')\)。故固定模生成元后，任意两组关系生成元给出相同理想。

	其次，向原来的生成元 \(m_1,\ldots,m_n\) 中添加一个冗余生成元
	\[
		m_{n+1}=\sum_{j=1}^n a_jm_j.
	\]
	把 \(a=(a_1,\ldots,a_n)^{\mathsf T}\) 看作列向量。新满射 \(A^{n+1}\to M\) 的关系模由矩阵
	\[
		G=\begin{pmatrix}H&-a\\0&1\end{pmatrix}
	\]
	的各列生成。确实，若 \((u,v)\in A^n\oplus A\) 是关系，就有 \(u+av\in\ker\pi=\operatorname{im}H\)，从而可写为
	\[
		(u,v)=(Hb-av,v)=G(b,v).
	\]
	反方向直接代入满射即可验证。左乘可逆矩阵消去右上块后，得到
	\[
		\begin{pmatrix}I_n&a\\0&1\end{pmatrix}G
		=\begin{pmatrix}H&0\\0&1\end{pmatrix}.
	\]
	于是对 \(0\leq i\leq n\)，令 \(q=n+1-i\geq1\)，上个引理给出
	\[
		I_{n+1-i}(G)=I_{n-i}(H).
	\]
	当 \(i=n\) 时右侧按 \(I_0(H)=A\) 理解；当 \(i>n\) 时，新旧定义都为 \(A\)。所以添加一个生成元及其定义关系不会改变 Fitting 理想。

	最后，给定两组有限生成元 \(m_1,\ldots,m_n\) 和 \(m'_1,\ldots,m'_{n'}\)，从第一组开始逐个添加第二组，或从第二组开始逐个添加第一组，最终得到同一组拼接生成元，至多相差排列。每次添加都是上面已处理的冗余生成元；排列是自由模的可逆换基，也由引理保持子式理想。对拼接后的固定生成元，第一步又保证任意关系生成组给出相同理想。这就比较了任意两个有限呈示。若更换为同构模，只需沿同构传送满射与生成元，相同证明仍适用。
\end{proof}

有限呈示保证整个计算只使用有限矩阵。证明没有把某个呈示当作唯一的呈示，也没有要求两组生成元有相同项数。下一项性质因此可以在任意方便的呈示上验证。

\begin{proposition}\label{g5:prop:fitting-base-change}
	设 \(A\to B\) 为任意交换含幺环同态，\(M\) 为有限呈示 \(A\) 模。则 \(M\otimes_A B\) 是有限呈示 \(B\) 模，且对所有整数 \(i\geq-1\)，有
	\[
		\operatorname{Fitt}_i(M\otimes_A B)
		=\operatorname{Fitt}_i(M)B.
	\]
	左侧理想在 \(B\) 中计算，右侧是扩张理想；这里不要求 \(B\) 对 \(A\) 平坦。
\end{proposition}
\begin{proof}
	对\eqref{g5:eq:fitting-presentation}张量 \(B\)。张量积的右正合性给出
	\[
		B^m\xrightarrow{H_B}B^n
		\longrightarrow M\otimes_A B\longrightarrow0,
	\]
	其中 \(H_B\) 由逐项变换 \(H\) 的系数得到。行列式是矩阵条目的整数系数多项式，所以 \(H_B\) 的每个子式恰为对应原子式在 \(B\) 中的像，生成的理想正是原子式理想的扩张。单位理想与零理想的约定也与基变换相容。这里只使用右正合性，没有要求张量后保留 \(A^m\) 左侧的任何单射。
\end{proof}

特别地，对乘法集 \(S\subseteq A\)，有
\[
	\operatorname{Fitt}_i(S^{-1}M)
	=S^{-1}\operatorname{Fitt}_i(M).
\]
因此在一点处检验 Fitting 理想，可以先计算整体子式再局部化；在商环上检验，则可以直接把子式的系数取模。

\subsection{剩余域上的秩条件}
\label{g5:subsec:0602:03}

记
\[
	M(\mathfrak p)=M\otimes_A\kappa(\mathfrak p),\qquad
	\nu_M(\mathfrak p)=\dim_{\kappa(\mathfrak p)}M(\mathfrak p).
\]
对有限呈示模，这些维数有限。第二章已经证明它们由关系矩阵的秩控制；Fitting 理想现在把同一计算变成不依赖呈示的闭集公式。

\begin{proposition}\label{g5:prop:fitting-fiber-rank}
	设 \(A\) 为交换含幺环，\(M\) 为有限呈示 \(A\) 模，\(\mathfrak p\in\operatorname{Spec}A\)。对整数 \(i\geq0\)，以下条件等价：
	\begin{equivalences}
		\item \(\nu_M(\mathfrak p)\leq i\)；
		\item \(\operatorname{Fitt}_i(M)\nsubseteq\mathfrak p\)；
		\item \(M_{\mathfrak p}\) 可由至多 \(i\) 个元素生成。
	\end{equivalences}
	特别地，对整数 \(r\geq0\)，有
	\[
		\{\mathfrak p:\nu_M(\mathfrak p)=r\}
		=V\bigl(\operatorname{Fitt}_{r-1}(M)\bigr)
		\setminus V\bigl(\operatorname{Fitt}_r(M)\bigr),
	\]
	而对 \(i\geq0\)，有
	\[
		\{\mathfrak p:\nu_M(\mathfrak p)>i\}
		=V\bigl(\operatorname{Fitt}_i(M)\bigr).
	\]
\end{proposition}
\begin{proof}
	将有限呈示张量到 \(\kappa(\mathfrak p)\)，得到
	\[
		\nu_M(\mathfrak p)
		=n-\operatorname{rank}_{\kappa(\mathfrak p)}H(\mathfrak p).
	\]
	当 \(0\leq i<n\) 时，\(\nu_M(\mathfrak p)\leq i\) 等价于 \(H(\mathfrak p)\) 至少有一个 \(n-i\) 阶子式非零。一个 \(A\) 中的元素在 \(\kappa(\mathfrak p)\) 中非零，当且仅当它不属于 \(\mathfrak p\)，故这恰是第二条件。当 \(i\geq n\) 时，前两个条件都自动成立。

	若 \(M_{\mathfrak p}\) 可由 \(i\) 个元素生成，模去极大理想后得到的向量空间也至多 \(i\) 维。反过来，提升纤维的一组基，由\cref{g5:cor:nakayama-generators}生成整个有限模 \(M_{\mathfrak p}\)。这证明第三条件与第一条件等价。

	取补集便得到 \(\{\nu_M>i\}\) 的公式。对 \(r\geq1\)，精确等于 \(r\) 是“至少为 \(r\)”与“至多为 \(r\)”的交，得到第一个集合公式。\(r=0\) 时 \(V(\operatorname{Fitt}_{-1}(M))=V(0)=\operatorname{Spec}A\)，同一公式仍成立。
\end{proof}

因此纤维维数为固定值的点集是一个闭集与一个开集的交；维数增大的条件由 Fitting 理想的零点集给出。这恢复了\cref{g5:prop:fiber-upper-semicontinuity}，并说明其中的子式闭集如何内在地依赖于模。

\begin{corollary}\label{g5:cor:fitting-support}
	设 \(A\) 为交换含幺环，\(M\) 为有限呈示 \(A\) 模。则
	\[
		\operatorname{Supp}_A M=V\bigl(\operatorname{Fitt}_0(M)\bigr),
		\qquad
		\sqrt{\operatorname{Fitt}_0(M)}=\sqrt{\operatorname{Ann}_A M}.
	\]
\end{corollary}
\begin{proof}
	有限模 \(M_{\mathfrak p}\) 为零，当且仅当其纤维为零：从纤维消失到模消失使用 Nakayama 引理。所以上个命题在 \(i=0\) 时给出支集公式。再与\cref{g5:prop:support-annihilator}的 \(\operatorname{Supp}_A M=V(\operatorname{Ann}_A M)\) 比较，并用根式等于包含该理想的全部素理想之交，便得到根式相等。
\end{proof}

\begin{example}\label{g5:exm:fitting-crossing-support}
	设 \(A=k[x,y]\)，\(M=A/(x)\oplus A/(y)\)。对角呈示矩阵 \(H=\operatorname{diag}(x,y)\) 给出
	\[
		\operatorname{Fitt}_0(M)=(xy),\qquad
		\operatorname{Fitt}_1(M)=(x,y),\qquad
		\operatorname{Fitt}_2(M)=A.
	\]
	故支集为两条坐标轴之并；在原点纤维维数为二，在坐标轴上其余素点处为一，离开两条轴则为零。第二个理想专门记录两个生成方向同时出现的位置。

	Fitting 理想也未必等于零化子。令 \(A=k[t]\)、\(N=A/(t)\oplus A/(t)\)，则
	\[
		\operatorname{Fitt}_0(N)=(t^2),\qquad
		\operatorname{Ann}_A N=(t).
	\]
	二者有相同根式和零点集，却是不同理想。只留下支集，会丢掉这部分行列式信息。
\end{example}

\subsection{局部自由秩的理想判据}
\label{g5:subsec:0602:04}

纤维至多 \(r\) 维，只说明局部模可用 \(r\) 个元素生成；这些元素之间仍可能有非零关系。若改用这 \(r\) 个生成元写呈示，\(\operatorname{Fitt}_{r-1}\) 就由关系矩阵的所有一阶子式，即全部条目生成。因此令这个理想等于零，恰好会消去全部关系。

\begin{theorem}\label{g5:thm:fitting-locally-free}
	设 \(A\) 为交换含幺环，\(M\) 为有限呈示 \(A\) 模，\(r\geq0\) 为整数。则 \(M\) 局部自由且在每一点的秩都为 \(r\)，当且仅当
	\[
		\operatorname{Fitt}_r(M)=A,
		\qquad\operatorname{Fitt}_{r-1}(M)=0.
	\]
	更具体地，对每个 \(\mathfrak p\in\operatorname{Spec}A\)，有
	\[
		M_{\mathfrak p}\cong A_{\mathfrak p}^r
		\quad\Longleftrightarrow\quad
		\begin{cases}
			\operatorname{Fitt}_r(M)A_{\mathfrak p}=A_{\mathfrak p},\\
			\operatorname{Fitt}_{r-1}(M)A_{\mathfrak p}=0.
		\end{cases}
	\]
\end{theorem}
\begin{proof}
	先在非零局部环 \((R,\mathfrak m)\) 上证明相应断言，记有限呈示模为 \(P\)。若 \(P\cong R^r\)，自由模的计算与呈示无关性立即给出两个理想等式。

	反过来，设 \(\operatorname{Fitt}_r(P)=R\)、\(\operatorname{Fitt}_{r-1}(P)=0\)。若 \(r=0\)，由\cref{g5:prop:fitting-fiber-rank}，\(P/\mathfrak mP=0\)，Nakayama 引理给出 \(P=0\)。若 \(r\geq1\)，同一命题给出
	\[
		\dim_{R/\mathfrak m}P/\mathfrak mP\leq r,
		\qquad\dim_{R/\mathfrak m}P/\mathfrak mP>r-1,
	\]
	因此该维数恰为 \(r\)。提升一组基得到满射 \(R^r\to P\)。有限呈示性及\cref{g5:prop:finite-presentation-kernel}保证核有限生成，从而得到新的有限呈示
	\[
		R^s\xrightarrow{C}R^r\longrightarrow P\longrightarrow0.
	\]
	呈示无关性现在十分关键：用新的 \(r\) 个生成元计算，仍有
	\[
		I_1(C)=\operatorname{Fitt}_{r-1}(P)=0.
	\]
	因此 \(C\) 的每个条目都为零，满射的核也为零，得到 \(P\cong R^r\)。

	对 \(R=A_{\mathfrak p}\)、\(P=M_{\mathfrak p}\) 应用局部断言，再用 Fitting 理想与局部化相容，就得到逐点的等价条件。若整体两个理想等式成立，则每个局部模都自由且秩为 \(r\)；\cref{g5:lem:local-basis-extension}把每个点处的基扩展到主开邻域，得到所需局部自由性。

	最后，若 \(M\) 局部自由且秩为 \(r\)，则 \(\operatorname{Fitt}_r(M)\) 在每一点的局部化都是单位理想，故它不可能包含在任何极大理想中，只能等于 \(A\)。理想 \(\operatorname{Fitt}_{r-1}(M)\) 在每一点的局部化都为零，由\cref{g5:prop:local-detection}，它本身为零。零环情形没有素点，模也为零，陈述按同一约定成立。
\end{proof}

不能将第二个等式换成 \(\operatorname{Fitt}_{r-1}(M)\subseteq\mathfrak p\)。包含于 \(\mathfrak p\) 只保证纤维维数至少为 \(r\)，而等于零要求关系在局部环中已经消失。下一小节的例子会使这两个条件出现明确差别。

\begin{proposition}\label{g5:prop:finitely-presented-flat}
	设 \(A\) 为交换含幺环，\(M\) 为有限呈示 \(A\) 模。以下条件等价：
	\begin{equivalences}
		\item \(M\) 是平坦 \(A\) 模；
		\item \(M\) 是有限生成投射 \(A\) 模；
		\item \(M\) 局部自由，允许不同连通部分有不同的秩；
		\item 对每个 \(\mathfrak p\in\operatorname{Spec}A\)，记 \(r=\nu_M(\mathfrak p)\)，则
		\[
			\operatorname{Fitt}_{r-1}(M)A_{\mathfrak p}=0.
		\]
	\end{equivalences}
\end{proposition}
\begin{proof}
	前三个条件的等价性已在\cref{g5:thm:finite-presented-flat}中完整证明；其中有限呈示用于保证极小满射的关系模有限生成，平坦性再借助交关系 \(K\cap IF=IK\) 和 Nakayama 引理消去这个核。

	对给定 \(\mathfrak p\)，因为 \(r\) 正好是纤维维数，\cref{g5:prop:fitting-fiber-rank}自动给出 \(\operatorname{Fitt}_r(M)A_{\mathfrak p}=A_{\mathfrak p}\)。因此上个定理说明，第四条件等价于每个 \(M_{\mathfrak p}\) 都自由。有限呈示的局部基延伸引理又将此条件与第三条件等同。
\end{proof}

这里假设的是 \(M\) 作为 \(A\) 模有限呈示。一个有限呈示 \(A\) 代数通常含有无穷多个 \(A\) 模生成方向，例如 \(A[t]\)，所以不能只因它有有限多个代数生成元与关系，就套用本命题的有限模结论。

\ProseParagraphBreak

还可以把定理直接用于任意基变换。若 \(A\to B\)，则 \(M\otimes_A B\) 局部自由且常秩为 \(r\)，当且仅当
\[
	\operatorname{Fitt}_r(M)B=B,
	\qquad\operatorname{Fitt}_{r-1}(M)B=0.
\]
这由\cref{g5:prop:fitting-base-change,g5:thm:fitting-locally-free}立即得到。第一项要求相关子式在新基底上局部成为单位，第二项要求较低下标的 Fitting 理想被真正杀掉。取商环与局部化时，这两个要求可以分别用方程和主开集明确实现。

\subsection{幂零基底上恒定纤维维数的反例}
\label{g5:subsec:0602:05}

先回看\cref{g5:exm:constant-fiber-not-flat}。令 \(A=k[\varepsilon]/(\varepsilon^2)\)、\(M=A/(\varepsilon)\)。其唯一纤维是一维向量空间，而呈示矩阵 \((\varepsilon)\) 给出
\[
	\operatorname{Fitt}_1(M)=A,
	\qquad\operatorname{Fitt}_0(M)=(\varepsilon)\ne0.
\]
第一个条件说明一个生成元足够；第二个理想中的非零元素准确记录了关系 \(\varepsilon\overline1=0\)。纤维中 \(\varepsilon\) 为零，这条关系看不见，但局部自由秩一要求整个关系理想为零。因此 Fitting 判据直接指出了非平坦性所在。

\begin{example}\label{g5:exm:fitting-nilpotent-line}
	设 \(k\) 为域，
	\[
		A=k[t,\varepsilon]/(\varepsilon^2),
		\qquad M=A/(\varepsilon)\cong k[t].
	\]
	每个素理想都含有 \(\varepsilon\)，并通过取商与 \(k[t]\) 的素理想一一对应。因此对所有 \(\mathfrak p\in\operatorname{Spec}A\)，都有
	\[
		M\otimes_A\kappa(\mathfrak p)\cong\kappa(\mathfrak p),
	\]
	纤维维数在整条线上恒为一，包含泛型点也不例外。

	然而有限呈示 \(A\xrightarrow{\varepsilon}A\to M\to0\) 给出
	\[
		\operatorname{Fitt}_0(M)=(\varepsilon),
		\qquad\operatorname{Fitt}_1(M)=A.
	\]
	而且 \(\varepsilon\) 在每个局部环 \(A_{\mathfrak p}\) 中都非零。为验证这一点，每个元素唯一写成 \(a(t)+\varepsilon b(t)\)，乘以 \(\varepsilon\) 后得到 \(a(t)\varepsilon\)，所以
	\[
		\operatorname{Ann}_A(\varepsilon)=(\varepsilon).
	\]
	若 \(\varepsilon/1=0\) 于 \(A_{\mathfrak p}\)，就存在 \(s\notin\mathfrak p\) 使 \(s\varepsilon=0\)，但此时 \(s\in(\varepsilon)\subseteq\mathfrak p\)，矛盾。因此 \(M\) 在任何点处都不自由，也不平坦。

	也可直接看平坦性失败的单射。令 \(J=(\varepsilon)\)，则 \(J\cong A/(\varepsilon)\) 且 \(J^2=0\)。将包含 \(J\hookrightarrow A\) 张量 \(M\)，得到
	\[
		J\otimes_A M\cong J/\varepsilon J=J
		\longrightarrow A\otimes_A M\cong M.
	\]
	这个映射为零而源非零。取商基变换 \(A\to A/(\varepsilon)\) 后，\(M\) 却变为秩一自由模；Fitting 基变换公式相应地把 \((\varepsilon)\) 变成零理想。
\end{example}

从理想语言看，恒定纤维维数 \(r\) 给出的整体信息是
\[
	\operatorname{Fitt}_r(M)=A,
	\qquad
	\operatorname{Fitt}_{r-1}(M)\subseteq\sqrt{(0)}.
\]
第一个等式来自该理想不包含于任何素理想；第二个包含来自每个素理想都包含它。两者与纤维秩公式等价。局部自由常秩 \(r\) 则把第二个包含加强为等于零。

\ProseParagraphBreak

这也解释约化基环的作用。若 \(A\) 约化，则 \(\sqrt{(0)}=0\)，所以有限呈示模的纤维维数恒为 \(r\) 已足以推出局部自由常秩 \(r\)。若纤维维数只局部常值，就在每个维数恒定的主开邻域上应用同一论证；局部化仍约化，故也得到局部自由。这以 Fitting 理想重新说明了第二章\cref{g5:ex:reduced-constant-fibre}中的结论，同时指出幂零基底反例的区别：剩余域中的所有秩计算都成立，未消失的却是原环中的幂零关系理想。
% 纲要条目排为编号小节标题，并进入目录与书签；说明段仍保留于纲要注释中。

% !TeX root = ../../Book5.tex
% 纲要标题：### 6.3 泛自由性的证明
% 纲要位置：几何学计划纲要.md，第 5240 行。
% 正文状态：初稿；尚未正式审核。

% 纲要内容开始
% 
% * Noether 整基环上的有限型代数
% * 分式域上的标准单项式
% * 有限关系与清除分母
% * 局部化后的自由性
% * 自由秩可能无限的说明
% 
% 纲要内容结束

\section{泛自由性的证明}
\label{g5:sec:0603}

设 \(A\) 为整环，分式域为 \(K\)。把一个 \(A\) 模 \(M\) 变成 \(M\otimes_A K\) 后，总能在这个向量空间中选择一组基。几何上，这相当于先研究基底泛型点上最容易处理的纤维。要把所得基延伸到一个主开集 \(D(s)\)，必须找到同一个非零元素 \(s\in A\)，使全部需要的系数和关系都在 \(A_s\) 上成立。

\ProseParagraphBreak

当 \(M\) 只在某个有限型 \(A\) 代数 \(B\) 上有限生成时，它作为 \(A\) 模可能需要无限多个生成元，分式域上的基也可能是无限集。本节用多项式关系控制这组基：通过有限个首一关系，统一消去所有不属于基的单项式。有限性落在关系及其首项上，因此只需清除有限多个分母。

\subsection{Noether 整基环上的有限型代数}
\label{g5:subsec:0603:01}

\begin{theorem}[泛自由性]\label{g5:thm:generic-freeness}
	设 \(A\) 为 Noether 整环，\(B\) 为有限型 \(A\) 代数，\(M\) 为有限生成 \(B\) 模。则存在 \(0\ne s\in A\) 及集合 \(\mathcal E\)，使
	\[
		M_s\cong\bigoplus_{e\in\mathcal E}A_s
	\]
	作为 \(A_s\) 模同构。这里 \(\mathcal E\) 可以为空集，也可以是无限集；结论不要求 \(A\to B\) 单射。
\end{theorem}

这就是\term[fanziyouxing]{泛自由性}[generic freeness]：在一个包含泛型点的非空主开集上，模具有统一的一组基。定理是 \cite[第24节，定理24.1，第185—186页]{Matsumura1989}中更一般结论的特例。下面给出标准单项式证明；所需的模版本除法及清分母论证分别建立后，在\subsecref{g5:subsec:0603:04}完成定理的证明。

\ProseParagraphBreak

有限型和有限生成这两个条件先把问题变成有限个多项式关系。取 \(B\) 的有限组 \(A\) 代数生成元，得到满同态
\[
	P=A[x_1,\ldots,x_n]\longrightarrow B.
\]
再取 \(M\) 的有限组 \(B\) 模生成元，就得到满同态 \(P^r\to M\)；允许加入零生成元，从而可取 \(r\geq1\)。由 Hilbert 基定理 \cref{g5:thm:hilbert-basis}逐次添加变量，\(P\) 为 Noether 环。于是\cref{g5:prop:noetherian-finite-modules}保证关系模
\[
	N=\ker(P^r\longrightarrow M)
\]
有限生成。固定生成元 \(h_1,\ldots,h_\ell\)，便有
\begin{equation}\label{g5:eq:generic-free-presentation}
	M\cong P^r/N,
	\qquad
	N=\sum_{i=1}^{\ell}Ph_i.
\end{equation}
零关系模允许 \(\ell=0\)。

\ProseParagraphBreak

令 \(K=\operatorname{Frac}A\)，并记
\[
	P_K=K[x_1,\ldots,x_n],\qquad
	N_K=K\otimes_A N\subseteq P_K^r.
\]
这里的包含以及
\[
	M\otimes_A K\cong P_K^r/N_K
\]
均来自局部化的正合性 \cref{g5:thm:localization-exact}。在 \(K\) 上，每个非零首项系数都可除去；我们将先利用这一点选择关系，再把所用的有限份数据移回某个 \(A_s\)。

\subsection{分式域上的标准单项式}
\label{g5:subsec:0603:02}

\cref{g5:prop:standard-monomials}已说明：对域上多项式理想，避开初始理想的单项式给出商环的向量空间基。现在关系是 \(r\) 个多项式组成的向量，因此要把各坐标的位置一起记入单项式。

\begin{definition}\label{g5:def:module-monomials}
	设 \(F\) 为域，\(n\geq0\)、\(r\geq1\) 为整数，\(Q=F[x_1,\ldots,x_n]\)，\(e_1,\ldots,e_r\) 为 \(Q^r\) 的标准基。形如 \(x^\alpha e_i\) 的向量称为\term[modanxiangshi]{模单项式}[module monomial]。规定
	\[
		x^\alpha e_i\mid x^\beta e_j
		\quad\Longleftrightarrow\quad
		i=j\ \text{且}\ \alpha_\nu\leq\beta_\nu\ (1\leq\nu\leq n).
	\]
	在 \(Q\) 中固定一个全局单项式序，并按下列规则比较模单项式：先比较 \(x^\alpha\)，相同时再按 \(e_1\prec\cdots\prec e_r\) 比较位置。对非零 \(v\in Q^r\)，记最大的模单项式、其系数及两者之积分别为
	\(\operatorname{LM}(v)\)、\(\operatorname{LC}(v)\) 和 \(\operatorname{LT}(v)\)。
\end{definition}

这个次序是良序。一个非空模单项式集合先有最小的多项式单项式，再在有限个可能的位置中取最小者；同时，它与乘以任意多项式单项式相容。因此，从向量中消去首项，只要引入的各项都更小，就不可能无限进行。

\ProseParagraphBreak

对 \(Q\) 子模 \(L\subseteq Q^r\)，定义其\term[chushizimo]{初始子模}[initial submodule]
\[
	\operatorname{in}(L)
	=\sum_{0\ne v\in L}Q\operatorname{LM}(v).
\]
这里仍用全部 \(L\) 中元素的首单项式。给定生成元的首项可能相消，从而留下新的首单项式；只取原生成元的首单项式，通常得不到 \(\operatorname{in}(L)\)。例如在 \(F[x,y]^2\) 中令
\[
	L=Q(xe_1+e_2)+Q(xe_1+ye_2),
\]
取字典序 \(x\succ y\)。两个生成元的首单项式均为 \(xe_1\)，而两者之差为 \((y-1)e_2\)，其首单项式 \(ye_2\) 不被 \(xe_1\) 整除。

\begin{lemma}\label{g5:lem:module-standard-monomials}
	设 \(F\) 为域，\(n\geq0\)、\(r\geq1\) 为整数，\(Q=F[x_1,\ldots,x_n]\)，\(L\subseteq Q^r\) 为子模，并采用上述模单项式序。则存在有限个首项系数为 \(1\) 的向量 \(g_1,\ldots,g_q\in L\)，使
	\[
		\operatorname{in}(L)
		=\sum_{j=1}^qQ\operatorname{LM}(g_j).
	\]
	这些 \(g_j\) 生成 \(L\)。所有不被任一 \(\operatorname{LM}(g_j)\) 整除的模单项式组成集合 \(\mathcal E\)，它们在 \(Q^r/L\) 中的剩余类构成一个 \(F\) 向量空间基。\(L=0\) 时允许 \(q=0\)。
\end{lemma}
\begin{proof}
	由 Hilbert 基定理及\cref{g5:prop:noetherian-finite-modules}，\(Q^r\) 的子模 \(\operatorname{in}(L)\) 有限生成。每个生成元都是有限个 \(\operatorname{LM}(v)\)、\(v\in L\) 的 \(Q\) 线性组合，因此可以从这些首单项式中选出有限个，仍生成整个初始子模。取相应的非零向量 \(g_j\in L\)，再除以首项系数，便可使它们首一。

	由模单项式生成的子模具有逐项判据：一个向量属于它，当且仅当其每个非零项的模单项式都被某个生成模单项式整除。一个方向由展开线性组合可知，因为每一项都处在相应坐标的单项式倍集中；反过来，逐项取倍式即可。特别地，\(0\ne v\in L\) 的首单项式一定被某个 \(\operatorname{LM}(g_j)\) 整除。

	对任意 \(v\in Q^r\) 实行除法。起初令待处理向量 \(p=v\)，余项 \(u=0\)。若
	\[
		\operatorname{LT}(p)=c x^\alpha e_i,
		\qquad
		\operatorname{LM}(g_j)=x^\beta e_i\mid x^\alpha e_i,
	\]
	就以
	\[
		p-cx^{\alpha-\beta}g_j
	\]
	替换 \(p\)，并记录所减的倍式。由于 \(g_j\) 首一，首项恰好消去，余下各项都更小。若没有这样的 \(j\)，就把 \(\operatorname{LT}(p)\) 移入 \(u\)，从 \(p\) 中删去此项。无论采取哪一步，非零待处理向量的首单项式都严格下降；良序性保证有限步后 \(p=0\)。从始至终记录的等式给出
	\[
		v=\sum_{j=1}^q a_jg_j+u,
		\qquad a_j\in Q,
	\]
	而 \(u\) 的每个模单项式都属于 \(\mathcal E\)。

	若 \(v\in L\)，则 \(u\in L\)。非零的 \(u\) 一方面全由 \(\mathcal E\) 中单项式组成，另一方面其首单项式必须被某个 \(\operatorname{LM}(g_j)\) 整除，矛盾。因此 \(u=0\)，也就证明了 \(g_j\) 生成 \(L\)。对一般 \(v\)，除法说明 \(\mathcal E\) 的类张成商空间。若这些类的一个非零有限线性组合属于 \(L\)，对其首单项式作同一判断又得到矛盾，故这些类线性无关。
\end{proof}

这些向量就是模的 Gröbner 基；所用的生成性与余式论证和\cref{g5:prop:groebner-remainder}相同，只多记录了坐标位置。这里的存在性来自初始子模的 Noether 性，随后逐项消去给出有限终止的除法。标准单项式集合 \(\mathcal E\) 本身无需有限。

\subsection{有限关系与清除分母}
\label{g5:subsec:0603:03}

将引理用于 \(P_K^r\) 的子模 \(N_K\)，得到首一关系 \(g_1,\ldots,g_q\) 及标准单项式集合 \(\mathcal E\)。现在须把同一组关系用于 \(A_s\)。仅使 \(g_j\) 的坐标属于 \(A_s[x_1,\ldots,x_n]\)，还没有说明它们属于原关系模 \(N_s\)；仅使它们属于 \(N_s\)，也尚未说明它们生成 \(N_s\)。两方向都可由有限份表达式落实。

\begin{lemma}\label{g5:lem:spread-monic-relations}
	设 \(A\) 为整环，\(K=\operatorname{Frac}A\)，\(n\geq0\)、\(r\geq1\) 为整数，并记 \(P=A[x_1,\ldots,x_n]\)。采用上述模单项式序，且设
	\[
		N=\sum_{i=1}^{\ell}Ph_i\subseteq P^r
	\]
	有限生成。设 \(g_1,\ldots,g_q\) 是 \(N_K\subseteq K[x_1,\ldots,x_n]^r\) 的有限个首一生成元。则存在 \(0\ne s\in A\)，使每个 \(g_j\) 都属于 \(N_s\)，并且
	\[
		N_s=\sum_{j=1}^q A_s[x_1,\ldots,x_n]g_j.
	\]
	这些向量在 \(A_s\) 上的首项系数仍为 \(1\)，首单项式也不变。
\end{lemma}
\begin{proof}
	在 \(P_K^r\) 内选择表达式
	\[
		g_j=\sum_{i=1}^{\ell}a_{ji}h_i,
		\qquad
		h_i=\sum_{j=1}^q b_{ij}g_j,
		\qquad a_{ji},b_{ij}\in K[x_1,\ldots,x_n].
	\]
	第一组等式来自 \(N_K\) 由 \(h_i\) 生成，第二组来自 \(g_j\) 的生成性。这里共有有限个多项式，每个又只有有限个系数。把它们所有系数写成 \(A\) 中元素的分式，令 \(s\) 为全部非零分母的乘积；没有分母时取 \(s=1\)。由于 \(A\) 为整环，\(s\ne0\)。每个分母均在 \(A_s\) 中可逆，因而全部 \(a_{ji},b_{ij}\) 都属于 \(A_s[x_1,\ldots,x_n]\)。

	第一组等式使 \(g_j\in N_s\)，同时也使 \(g_j\) 的全部坐标属于 \(A_s[x_1,\ldots,x_n]\)。第二组等式使每个 \(h_i\) 都属于这些 \(g_j\) 生成的子模，故两边相等。包含 \(A_s\hookrightarrow K\) 保持非零系数及全部多项式等式，所以同一个 \(g_j\) 的首单项式和首项系数均不变。若 \(N_K=0\)，则 \(P^r\hookrightarrow P_K^r\) 给出 \(N=0\)，此时取空生成元表和 \(s=1\) 即可。
\end{proof}

此处清除的是有限个关系及相互表示的系数。标准单项式即使有无限多个，也没有参加取分母的乘积。只从 \(M\otimes_A K\) 的任意一组无限基出发逐个选择提升，可能遇到无限多分母，因而不能直接完成同一个主开集上的构造。

\subsection{局部化后的自由性}
\label{g5:subsec:0603:04}

选择上个引理给出的 \(s\)，记 \(P_s=A_s[x_1,\ldots,x_n]\)。现在 \(g_j\in N_s\subseteq P_s^r\)，且每个首项系数都是单位 \(1\)。虽然 \(A_s\) 一般仍不是域，模除法中实际需要的除法已经全部可做。

\begin{proof}[\cref{g5:thm:generic-freeness}的证明]
	采用\eqref{g5:eq:generic-free-presentation}中的 \(P^r/N\)。由\cref{g5:lem:module-standard-monomials}为 \(N_K\) 选择首一关系 \(g_1,\ldots,g_q\)，再由\cref{g5:lem:spread-monic-relations}选择 \(0\ne s\in A\)，使这些关系生成 \(N_s\)。令 \(\mathcal E\) 为不被其任何首单项式整除的模单项式集合。

	局部化正合性给出 \(M_s\cong P_s^r/N_s\)。对 \(v\in P_s^r\)，逐步作\cref{g5:lem:module-standard-monomials}证明中的除法。每次消去 \(cx^\alpha e_i\) 时，只需减去 \(cx^{\alpha-\beta}g_j\)，系数一直留在 \(A_s\) 中；移动不能消去的项也不改变系数环。相同的良序论证保证终止。因此 \(v\) 模 \(N_s\) 是 \(\mathcal E\) 中有限多个单项式的 \(A_s\) 线性组合，得到满同态
	\[
		\Phi:\bigoplus_{e\in\mathcal E}A_s\longrightarrow M_s,
		\qquad (c_e)_e\longmapsto\sum_e c_e[e].
	\]

	若某个非零有限组合 \(w=\sum_e c_e e\) 落在 \(N_s\)，由 \(A_s\hookrightarrow K\)，它在 \(P_K^r\) 中仍非零，且属于 \(N_K\)。于是其首单项式被某个 \(\operatorname{LM}(g_j)\) 整除。但 \(w\) 的全部单项式都在 \(\mathcal E\) 中，矛盾。因此 \(\Phi\) 也单射，这就得到所需自由基。\(\mathcal E=\varnothing\) 时结论是 \(M_s=0\)，即秩为零的自由模。
\end{proof}

这组基同时是泛型纤维的基：再次局部化便有
\[
	M\otimes_A K\cong\bigoplus_{e\in\mathcal E}K.
\]
更一般地，对每个 \(\mathfrak p\in D(s)\)，张量积与直和相容，给出
\[
	M\otimes_A\kappa(\mathfrak p)
	\cong\bigoplus_{e\in\mathcal E}\kappa(\mathfrak p).
\]
所以 \(D(s)\) 上各模纤维都有同一个基的指标集。当该集合无限时，这里说的是无限维向量空间。若 \(M\) 本身有限生成于 \(A\)，则其泛型纤维由有限多个向量生成，所以 \(\mathcal E\) 必定有限。

\begin{corollary}\label{g5:cor:simultaneous-generic-freeness}
	设 \(A\) 为 Noether 整环，\(B\) 为有限型 \(A\) 代数，\(M_1,\ldots,M_m\) 为有限多个有限生成 \(B\) 模。存在同一个 \(0\ne s\in A\)，使每个 \((M_i)_s\) 都是自由 \(A_s\) 模。还可同时要求 \(B_s\) 自由。
\end{corollary}
\begin{proof}
	分别应用泛自由性定理，选择非零元素 \(s_i\in A\)，使 \((M_i)_{s_i}\) 自由。需要时也把有限生成 \(B\) 模 \(B\) 列入这张有限表。取所有这些元素的乘积 \(s\)，则 \(s\ne0\)。每个 \(s_i\) 在 \(A_s\) 中可逆，故 \((M_i)_s\) 是 \((M_i)_{s_i}\) 的进一步局部化。局部化保持直和及其标准基，所以仍自由。
\end{proof}

自由模即使秩无限也平坦：与 \(\bigoplus_{e\in\mathcal E}A_s\) 张量，等于取 \(\mathcal E\) 份原映射的直和，而直和逐坐标保持单射。因此泛自由性给出了同一个主开集上的泛平坦性。\secref{g5:sec:0604}将把这一结论用于几何族，并区分泛型开集与全部平坦位置。

\subsection{自由秩可能无限的说明}
\label{g5:subsec:0603:05}

\begin{example}\label{g5:exm:generic-free-infinite-rank}
	设 \(A\) 为 Noether 整环，\(B=A[x]\)，\(M=B\)。\(M\) 是由 \(1\) 生成的 \(B\) 模，却作为 \(A\) 模具有基
	\[
		1,x,x^2,x^3,\ldots.
	\]
	多项式唯一的系数展开同时证明张成性与线性无关性。对任意 \(0\ne s\in A\)，这仍是 \(A_s[x]\) 的 \(A_s\) 基；再到分式域得到无限维空间 \(K[x]\)，因此不可能通过缩小主开集把它变成有限秩自由模。
\end{example}

有限型条件控制作为代数所需的变量数，而有限生成 \(B\) 模条件控制相对于这些变量所需的模生成元数。变量的任意次幂仍可提供无限多个 \(A\) 线性无关的元素。下个例子同时展示关系中的一个系数如何决定应删去的基底位置。

\begin{example}\label{g5:exm:generic-free-tx}
	设 \(k\) 为域，\(A=k[t]\)，\(B=A[x]\)，并取有限生成 \(B\) 模
	\[
		M=B/(tx).
	\]
	理想 \((tx)\) 中的多项式常数项为零，其每个正次幂系数均被 \(t\) 整除；反过来，满足这两个条件的多项式都属于 \((tx)\)。按系数逐项比较，得到 \(A\) 模分解
	\[
		M\cong A\oplus\bigoplus_{j\geq1}(A/(t))x^j.
	\]
	特别地，\([x]\ne0\) 而 \(t[x]=0\)。把单射 \(A\xrightarrow{\,t\,}A\) 与 \(M\) 张量后，就得到不是单射的 \(M\xrightarrow{\,t\,}M\)，所以 \(M\) 在 \(A\) 上不平坦。

	在 \(K=k(t)\) 上，关系 \(tx\) 可以除以 \(t\)，变为首一关系 \(x\)。标准单项式只有 \(1\)。两方向的关系表达式是
	\[
		x=t^{-1}(tx),\qquad tx=t\cdot x,
	\]
	所以清除分母只需选择 \(s=t\)，并且
	\[
		M_t\cong A_t[x]/(x)\cong A_t.
	\]
	在 \(\mathfrak p=(t)\) 处，模纤维是 \(k[x]\)，维数无限；在任一 \(\mathfrak p\in D(t)\) 处，它则是 \(\kappa(\mathfrak p)\)，维数为 \(1\)。退化的原因正是关系的系数 \(t\) 在特殊点变成了零。
\end{example}

\begin{example}\label{g5:exm:generic-free-no-common-open}
	设 \(k\) 为无限域，\(A=B=k[t]\)，对每个 \(a\in k\) 取有限生成模 \(M_a=A/(t-a)\)。分别反演 \(t-a\)，可使 \((M_a)_{t-a}=0\)，从而自由。但不存在适用于所有 \(M_a\) 的同一个非零 \(s\in A\)。事实上，非零多项式 \(s\) 只有有限多个根，可取 \(a\in k\) 使 \(s(a)\ne0\)。于是
	\[
		(M_a)_s\cong A_s/(t-a)A_s\cong k\ne0.
	\]
	它被 \(A_s\) 中的非零元素 \(t-a\) 零化，故不是整环 \(A_s\) 上的自由模。这说明同时选择一个主开集时，模的数目有限这一条件确有作用。
\end{example}

泛自由性把基底上的系数变化归结为有限个分母：首项系数在 \(D(s)\) 上可逆后，同一套消去规则适用于全部多项式向量，也适用于这个主开集上的每个模纤维。它给出的统一线性结构允许无限秩，因而适合处理有限型代数及其有限模，而不限于基底上的有限模。
% 纲要条目排为编号小节标题，并进入目录与书签；说明段仍保留于纲要注释中。

% !TeX root = ../../Book5.tex
% 纲要标题：### 6.4 忠实平坦与几何族
% 纲要位置：几何学计划纲要.md，第 5248 行。
% 正文状态：初稿；尚未正式审核。

% 纲要内容开始
% 
% * 消失和正合性的反映
% * 谱映射满射的判据
% * 泛平坦性及开放条件
% * 平坦退化的显式方程
% * 为下降准备的双重基变换
% 
% ---
% 
% 纲要内容结束

\section{忠实平坦与几何族}
\label{g5:sec:0604}

平坦的系数扩张保持正合列，却可能把非零对象整个消去。例如把整数扩张到有理数后，每个有限阿贝尔群都变为零。若希望在新系数下验证一个结论，再把结论传回原环，就需要保证这种消失不会发生。忠实平坦性正是平坦性与检测能力的结合。先建立这一代数条件，再将它用于纤维、泛型开集和可显式计算的退化。

\subsection{消失和正合性的反映}
\label{g5:subsec:0604:01}

\begin{definition}\label{g5:def:faithfully-flat-algebra}
	设 \(A\to B\) 为交换含幺环同态。若 \(B\) 是平坦 \(A\) 模，且对任意 \(A\) 模 \(N\) 都有
	\[
		B\otimes_A N=0\quad\Longrightarrow\quad N=0,
	\]
	则称 \(A\to B\) 为\term[zhongshipingtan]{忠实平坦}[faithfully flat]同态，也称 \(B\) 为忠实平坦 \(A\) 代数。
\end{definition}

这里的“忠实”要求张量函子不丢掉任何非零模。它强于环同态单射：\(\mathbb Z\hookrightarrow\mathbb Q\) 虽然单射且平坦，却把 \(\mathbb Z/2\mathbb Z\) 张量成零。

\begin{proposition}\label{g5:prop:faithful-flat-reflects-exactness}
	设 \(A\to B\) 为忠实平坦同态。一个 \(A\) 模序列
	\[
		L\xrightarrow{u}M\xrightarrow{v}N
	\]
	在中间项正合，当且仅当张量 \(B\) 后在中间项正合。对任意 \(A\) 模同态，单射、满射、同构以及零映射也都可在张量 \(B\) 后检测。
\end{proposition}
\begin{proof}
	平坦性保持短正合列，因此保持任意同态的核与像：把同态分解为到其像的满射和像的包含，分别应用短正合列即可。它也保持余核，这是张量的右正合性。

	若张量后的序列正合，则 \((vu)\otimes1=0\)，故
	\[
		B\otimes_A\operatorname{im}(vu)
		\cong\operatorname{im}\bigl((vu)\otimes1\bigr)=0.
	\]
	忠实性使 \(\operatorname{im}(vu)=0\)，即 \(vu=0\)。于是 \(\operatorname{im}u\subseteq\ker v\)，可以形成商模
	\(H=\ker v/\operatorname{im}u\)。平坦性给出
	\[
		B\otimes_A H
		\cong\ker(v\otimes1)/\operatorname{im}(u\otimes1)=0.
	\]
	再次用忠实性得 \(H=0\)，所求正合性成立。反方向只用平坦性。把同样的消失检测分别用于核、余核、像，就得到其余断言。
\end{proof}

因此，在忠实平坦扩张后核验一个矩阵关系或一列模的正合性，可以传回原环；平坦性负责正确运送核与像，忠实性负责从运送后的消失恢复原来的消失。这个等价刻画可对照 \cite[第7节，定理7.2，第47页]{Matsumura1989}。

\begin{proposition}\label{g5:prop:faithful-flat-base-change}
	忠实平坦同态在复合及任意基变换下保持成立。具体地，若 \(A\to B\) 忠实平坦、\(A'\) 为任意 \(A\) 代数，则 \(A'\to B\otimes_A A'\) 忠实平坦；若 \(A\to B\to C\) 的两步都忠实平坦，则 \(A\to C\) 忠实平坦。
\end{proposition}
\begin{proof}
	对任意 \(A'\) 模 \(N\)，有自然同构
	\[
		(B\otimes_A A')\otimes_{A'}N\cong B\otimes_A N.
	\]
	一个 \(A'\) 模的单射也是 \(A\) 模单射，一个非零 \(A'\) 模作为 \(A\) 模仍非零。因此右侧既保持单射，又不消去非零模，给出基变换断言。对复合使用
	\(C\otimes_A N\cong C\otimes_B(B\otimes_A N)\)：保持单射逐次成立，消失也可逐次反映。
\end{proof}

\subsection{谱映射满射的判据}
\label{g5:subsec:0604:02}

能否检测所有模，可以先从最简单的循环模 \(A/I\) 检查。它张量后变为 \(B/IB\)；若某个真理想扩张成整个 \(B\)，相应商模就已经消失。只需检查极大理想，就能排除全部这种情况。

\begin{theorem}\label{g5:thm:faithful-flat-spectrum}
	设 \(A\to B\) 为平坦的交换含幺环同态。以下条件等价：
	\begin{equivalences}
		\item \(A\to B\) 忠实平坦；
		\item 每个极大理想 \(\mathfrak m\subset A\) 都满足 \(\mathfrak mB\ne B\)；
		\item 每个真理想 \(I\subset A\) 都满足 \(IB\ne B\)；
		\item 取收缩的谱映射 \(\operatorname{Spec}B\to\operatorname{Spec}A\) 满射。
	\end{equivalences}
\end{theorem}
\begin{proof}
	忠实性用于非零模 \(A/\mathfrak m\)，立即给出第二条件。每个真理想都包含在某个极大理想内，故第二条件蕴含第三条件。

	设第三条件成立，取非零 \(A\) 模 \(N\) 中的非零元素 \(n\)。循环子模
	\(An\cong A/\operatorname{Ann}_A(n)\)，其零化子为真理想，所以
	\(B\otimes_A An\ne0\)。平坦性使 \(B\otimes_A An\to B\otimes_A N\) 单射，故后者也非零。这就恢复了忠实性。

	若忠实平坦，对每个 \(\mathfrak p\in\operatorname{Spec}A\)，非零 \(A\) 模 \(\kappa(\mathfrak p)\) 给出非零环
	\[
		B\otimes_A\kappa(\mathfrak p)
		\cong (A\setminus\mathfrak p)^{-1}(B/\mathfrak pB).
	\]
	非零含幺环存在素理想，商环与局部化的素理想对应将它送回某个 \(\mathfrak q\subset B\)，满足 \(\mathfrak q\cap A=\mathfrak p\)。这证明谱映射满射。反过来，若谱映射满射，对任意极大理想 \(\mathfrak m\) 取其上方素理想 \(\mathfrak q\)，则 \(\mathfrak mB\subseteq\mathfrak q\ne B\)，得到第二条件。零环情形各条件也按空谱约定成立。
\end{proof}

谱纤维在这里是纤维环的素谱，点可以具有扩大的剩余域。满射要求每个基底素点有某个上方素点，并不要求上方存在具有相同剩余域的有理点。上述判据及纤维环说明见 \cite[第7节，定理7.2—7.3，第47—48页]{Matsumura1989}。

\begin{corollary}\label{g5:cor:flat-local-faithful}
	设 \((A,\mathfrak m)\to(B,\mathfrak n)\) 为非零局部环之间的局部同态。若 \(B\) 对 \(A\) 平坦，则它忠实平坦。
\end{corollary}
\begin{proof}
	局部同态给出 \(\mathfrak mB\subseteq\mathfrak n\ne B\)，而 \(A\) 只有一个极大理想，故可用定理。
\end{proof}

一个局部化 \(A\to S^{-1}A\) 总是平坦，但通常不忠实。其谱像是与 \(S\) 不相交的素理想。这个像为整个谱，当且仅当 \(S\) 的每个元素都是单位：非单位属于某个极大理想，而单位不属于任何素理想。所以一个真正删去点的局部化也会失去对某些模的检测。

\begin{example}\label{g5:exm:faithful-principal-cover}
	若 \(f_1,\ldots,f_r\in A\)、\(r\geq1\)，且 \(D(f_i)\) 覆盖 \(\operatorname{Spec}A\)，则
	\[
		A\longrightarrow\prod_{i=1}^r A_{f_i}
	\]
	忠实平坦。有限乘积作为模就是有限直和，故平坦。每个极大理想 \(\mathfrak m\) 都有某个 \(f_i\notin\mathfrak m\)，于是
	\((A/\mathfrak m)_{f_i}\cong A/\mathfrak m\ne0\)，说明乘积模不被 \(\mathfrak m\) 全部生成。由定理得忠实性。这里把一组局部检查合并成了一次忠实平坦扩张。
\end{example}

\subsection{泛平坦性及开放条件}
\label{g5:subsec:0604:03}

泛自由性给出的 \(s\ne0\) 把坏行为限制在 \(V(s)\) 中。在 Noether 整基环上，\(D(s)\) 是包含泛型点的稠密开集：两个非空主开集的交仍含泛型点，故都与它相交。在这个开集上，自由性立即给出平坦性。

\begin{corollary}[泛平坦性]\label{g5:cor:generic-flatness}
	设 \(A\) 为 Noether 整环，\(B\) 为有限型 \(A\) 代数，\(M_1,\ldots,M_r\) 为有限生成 \(B\) 模。存在 \(0\ne s\in A\)，使 \(B_s\) 及所有 \((M_i)_s\) 都是平坦 \(A_s\) 模。可以进一步使它们都自由，允许无限秩和零模。
\end{corollary}
\begin{proof}
	对有限多个 \(B\) 模 \(B,M_1,\ldots,M_r\) 应用\cref{g5:cor:simultaneous-generic-freeness}，得到同一个 \(s\ne0\)，使各模在 \(A_s\) 上自由。任意秩的自由模都平坦，因而得到全部断言。
\end{proof}

“有限生成 \(B\) 模”与“有限呈示 \(A\) 模”在这里承担不同任务。前者允许 \(M\) 在 \(A\) 上有无限秩，只保证能找到一块泛型开集。对后者，则可以证明其全部平坦位置组成开集。

\begin{proposition}\label{g5:prop:finite-module-flat-locus}
	设 \(A\) 为交换含幺环，\(M\) 为有限呈示 \(A\) 模。集合
	\[
		U=\{\mathfrak p\in\operatorname{Spec}A:
		M_{\mathfrak p}\text{ 在 }A_{\mathfrak p}\text{ 上平坦}\}
	\]
	为开集。在 \(U\) 内，纤维维数局部常值。
\end{proposition}
\begin{proof}
	若 \(\mathfrak p\in U\)，则 \(M_{\mathfrak p}\) 仍有限呈示，由\cref{g5:thm:finite-presented-flat}在局部环上自由，设秩为 \(r\)。第二章的局部基延伸引理\cref{g5:lem:local-basis-extension}把这组基延伸到某个 \(D(f)\)、\(f\notin\mathfrak p\)，使 \(M_f\cong A_f^r\)。该主开集上的每个局部化都自由，因而平坦，所以 \(D(f)\subseteq U\)。在这个邻域上纤维维数也恒为 \(r\)。
\end{proof}

Fitting 理想把这个判定变成方程。对于固定的秩 \(r\)，在一点附近自由要求
\[
	\operatorname{Fitt}_r(M)_{\mathfrak p}=A_{\mathfrak p},
	\qquad
	\operatorname{Fitt}_{r-1}(M)_{\mathfrak p}=0.
\]
前一个条件要求某个子式在该点可逆，后一个条件要求另一组子式在整个局部环中为零，见\cref{g5:thm:fitting-locally-free}。只把后一组子式代入剩余域，得到的仅是它们属于 \(\mathfrak p\)，尚不足以让局部环中的关系消失。有限呈示保证这些理想有限生成，因此局部的消失也能在某个主开邻域保持：逐个生成元清去分母，再取共同乘积即可。

\subsection{平坦退化的显式方程}
\label{g5:subsec:0604:04}

下面把参数 \(t\) 看作基底坐标。要证明一个方程族对 \(k[t]\) 平坦，可以直接寻找一组在所有参数值下共用的模基。这个基不必有限，也不要求所有纤维有相同的分支数。

\begin{example}\label{g5:exm:flat-node-family}
	设 \(k\) 为域，\(A=k[t]\)，
	\[
		B=k[t,x,y]/(xy-t).
	\]
	则 \(B\) 为忠实平坦 \(A\) 代数，且作为 \(A\) 模有无限基
	\[
		1,\quad x,x^2,x^3,\ldots,\quad y,y^2,y^3,\ldots.
	\]
	把每个混合单项式中的 \(xy\) 换成 \(t\)，就证明这组元素张成。为了证明独立，先由消去 \(t\) 得到 \(B\cong k[x,y]\)，其中 \(t\) 作用为乘 \(xy\)。一个有限的 \(A\) 线性关系展开后，来自 \(t^j x^m\)、\(m\geq1\)，的项具有指数差 \(m>0\)，来自 \(t^j y^n\)、\(n\geq1\)，的项具有指数差 \(-n<0\)，来自 \(t^j\) 的项具有指数差零。每一类内部的指数对也彼此不同，故所有系数为零。这证明自由性；基非空使与任何非零模张量后仍非零，所以忠实平坦。

	在参数 \(a\in k\) 处，
	\[
		B\otimes_A k[t]/(t-a)\cong k[x,y]/(xy-a).
	\]
	当 \(a\ne0\) 时，\(x\) 为单位，纤维环为 \(k[x,x^{-1}]\)；当 \(a=0\) 时，纤维是两条坐标轴的并。它们都是一维，但零纤维在原点的切空间为二维，而非零纤维在有理点处的局部环正则。平坦性容许这样的奇点与分支变化。
\end{example}

\begin{example}\label{g5:exm:flat-two-points-collision}
	设 \(k\) 为特征不等于二的域，\(A=k[t]\)，取
	\[
		C=A[z]/(z^2-t).
	\]
	除以首一多项式给出唯一余式 \(a(t)+b(t)z\)，故 \(C\) 为秩二的自由 \(A\) 模。若 \(a\ne0\)，纤维 \(k[z]/(z^2-a)\) 在代数闭包上分成两个不同点；在 \(a=0\) 时，纤维为一个带幂零元的加厚点 \(k[z]/(z^2)\)。每个纤维作为 \(k\) 向量空间仍有维数二。在特征二时自由性不变，但非零参数纤维在代数闭包上也为二重根，所以不能沿用“两点碰合”的解释。
\end{example}

与这两个例子比较，取
\[
	D=k[t,z]/(z^2,tz).
\]
它作为 \(k[t]\) 模分解为 \(k[t]\oplus k z\)，其中 \(t z=0\)、\(z\ne0\)。非零 \(t\) 是整基环中的非零因子，却杀死了非零模元素，故 \(D\) 不平坦。非零参数的纤维为 \(k\)，零参数的纤维为 \(k[z]/(z^2)\)。所有纤维的 Krull 维数均为零，改变的是长度和参数方向上的关系。只比较纤维维数，会遗漏这项障碍。

\subsection{为下降准备的双重基变换}
\label{g5:subsec:0604:05}

一次基变换把 \(M\) 变为 \(B\otimes_A M\)。若一个元素确实来自 \(M\)，把它再沿两张映射
\[
	p_1,p_2:B\longrightarrow B\otimes_A B,\qquad
	p_1(b)=b\otimes1,\quad p_2(b)=1\otimes b
\]
分别拉回，所得元素应相同。忠实平坦性保证这个相容条件也足以识别来自原模的元素。

\begin{theorem}\label{g5:thm:faithful-flat-equalizer}
	设 \(A\to B\) 为忠实平坦的交换含幺环同态，\(M\) 为任意 \(A\) 模。以下 \(A\) 模序列正合：
	\[
		0\longrightarrow M\xrightarrow{\eta}
		B\otimes_A M\xrightarrow{\delta}
		B\otimes_A B\otimes_A M,
	\]
	其中
	\[
		\eta(m)=1\otimes m,\qquad
		\delta(b\otimes m)=1\otimes b\otimes m-b\otimes1\otimes m.
	\]
	因此两种拉回相等的元素，都唯一来自 \(M\)。
\end{theorem}
\begin{proof}
	由\cref{g5:prop:faithful-flat-reflects-exactness}，只需在左边再张量一个 \(B\)，验证所得序列正合。它的前两张映射为
	\[
		\begin{aligned}
			\eta'(b\otimes m)&=b\otimes1\otimes m,\\
			\delta'(b\otimes c\otimes m)
			&=b\otimes1\otimes c\otimes m
			-b\otimes c\otimes1\otimes m.
		\end{aligned}
	\]
	乘法映射 \(\mu(b\otimes c\otimes m)=bc\otimes m\) 满足 \(\mu\eta'=1\)，所以 \(\eta'\) 单射。又 \(\delta'\eta'=0\)。

	若 \(z\in\ker\delta'\)，对等式 \(\delta'(z)=0\) 应用“乘起前两个 \(B\) 因子”的映射
	\[
		h:B\otimes_A B\otimes_A B\otimes_A M
		\longrightarrow B\otimes_A B\otimes_A M,
		\quad
		h(b\otimes c\otimes d\otimes m)=bc\otimes d\otimes m,
	\]
	就得到 \(z-\eta'(\mu z)=0\)。故 \(z\in\operatorname{im}\eta'\)，证明张量后的中间正合性。忠实平坦性将单射与正合性传回原序列。
\end{proof}

\begin{corollary}\label{g5:cor:faithful-flat-contraction}
	若 \(A\to B\) 忠实平坦，则 \(A\to B\) 单射，且每个理想 \(I\subseteq A\) 都满足
	\[
		\{a\in A:a\text{ 在 }B\text{ 中的像属于 }IB\}=I.
	\]
\end{corollary}
\begin{proof}
	在定理中取 \(M=A\) 得单射；取 \(M=A/I\) 得
	\(A/I\hookrightarrow B/IB\)，其零元判据就是所述收缩公式。
\end{proof}

最后的单射与收缩公式也见 \cite[第7节，定理7.5，第49页]{Matsumura1989}。平坦局部化未必满足收缩等于原理想，因为它可能删去理想所支撑的部分；忠实性在这里确保了检测不丢失。

\ProseParagraphBreak

两张 \(p_i\) 都是 \(A\to B\) 的基变换，故都忠实平坦。在三重张量 \(B\otimes_A B\otimes_A B\) 上，还可以比较从第一份到第二份、第二份到第三份以及直接从第一份到第三份的识别。以后为一个 \(B\) 模指定下降资料时，便要求这些识别的复合相容。本章的等化子序列已经证明了元素层面的识别：先扩张系数，再在双重基变换中检查两种像相等，就能恢复原模中的唯一元素。


\begin{chaptersummary}
	平坦性要求张量不增加单射左端的核。有限关系判据把它表达为关系可由系数之间的关系生成；有限生成理想判据则把检验对象缩减为理想的包含。局部检测允许在素点处验证这些条件。整环上的平坦模必无挠，但一般无挠模仍可能在多个系数的关系中出现障碍；主理想整环上的理想判据使二者等价。

	\ProseParagraphBreak

	Fitting 理想保留呈示矩阵子式所生成的理想，与生成元和关系的选择无关，并与任意基变换相容。纤维维数由子式在剩余域中的秩决定，局部自由性还要求较大子式在局部环中真正消失。幂零基底上，所有纤维都可能看不见这些非零关系。

	\ProseParagraphBreak

	泛自由性先在分式域上选择有限个首项控制关系，再清除这些关系及原关系生成元之间的有限多个分母。留下的标准单项式在一个共同的主开集上成为自由基；它们可以无限多个。忠实平坦扩张进一步反映模的消失与正合性，其谱映射覆盖全部基底素点。双重基变换中的等化子序列则把这种检测落实为元素的唯一恢复。
\end{chaptersummary}

\BookExercises

\begin{exercise}\label{g5:ex:flat-cyclic-idempotent}
	设 \(A\) 为交换含幺环，\(I\subseteq A\) 为有限生成理想。证明 \(A/I\) 平坦，当且仅当 \(I=Ae\)，其中 \(e^2=e\)。说明为什么有限生成条件在证明中有作用。
\end{exercise}
\begin{solution}
	若 \(A/I\) 平坦，理想包含张量后给出单射
	\(I\otimes_A A/I\to A/I\)。它的源同构于 \(I/I^2\)，映射却为零，故 \(I=I^2\)。
	选生成元 \(x_1,\ldots,x_n\)，写 \(x_i=\sum_j a_{ij}x_j\)，其中全部 \(a_{ij}\in I\)。令 \(C=(a_{ij})\)，伴随矩阵恒等式使 \(\det(1-C)I=0\)。展开行列式可写成 \(\det(1-C)=1-e\)，其中 \(e\in I\)。故 \((1-e)I=0\)，每个 \(x\in I\) 满足 \(x=ex\)，所以 \(I=Ae\)。又 \(e\in I\)，故 \((1-e)e=0\)，即 \(e^2=e\)。

	反过来，若 \(I=Ae\)、\(e^2=e\)，则 \(A=Ae\oplus A(1-e)\)，且 \(A/I\cong A(1-e)\)，所以商模为自由模 \(A\) 的直和分量，因而平坦。有限生成性使等式 \(I=I^2\) 能用一个有限矩阵表达，再用行列式找到共同的 \(e\)；本题并未把这一有限矩阵步骤推广到任意理想。
\end{solution}

\begin{exercise}\label{g5:ex:normalization-not-flat}
	设 \(A\to B\) 平坦，证明对任意两个理想 \(I,J\subseteq A\)，有
	\((I\cap J)B=IB\cap JB\)。据此证明尖点的正规化
	\(A=k[t^2,t^3]\subseteq B=k[t]\) 不平坦。此模型可对照 \cite[第7节，定理7.4后的例，第49页]{Matsumura1989}。
\end{exercise}
\begin{solution}
	映射 \(A\to A/I\oplus A/J\)、\(a\mapsto(a+I,a+J)\) 的核为 \(I\cap J\)。张量平坦模 \(B\) 后仍保留这个核，右端变为 \(B/IB\oplus B/JB\)，其对应映射的核为 \(IB\cap JB\)。同时 \(I\cap J\) 张量后在 \(B\) 中的像为 \((I\cap J)B\)，得所需等式。

	在尖点环中取 \(I=t^2A\)、\(J=t^3A\)。\(A\) 的单项式指数为 \(0,2,3,4,\ldots\)，因此 \(I\cap J\) 恰由次数至少五的单项式张成，也就是 \((t^5,t^6)A\)。于是
	\[
		(I\cap J)B=t^5B,\qquad IB\cap JB=t^2B\cap t^3B=t^3B.
	\]
	二者不同，故正规化不平坦。它虽为有限的单射并在素谱上满射，仍不能成为忠实平坦扩张。
\end{solution}

\begin{exercise}\label{g5:ex:localization-flat-not-free}
	证明 \(\mathbb Z[1/2]\) 是无挠平坦 \(\mathbb Z\) 模，却不是自由模，也不是忠实平坦模。
\end{exercise}
\begin{solution}
	它是 \(\mathbb Q\) 的子模，故无挠；局部化的正合性给出平坦性。乘二在 \(\mathbb Z[1/2]\) 上满射，而在任何非零自由 \(\mathbb Z\) 模上都不满射，因为一个基向量不能是两倍的整数线性组合。因此它不自由。又
	\[
		\mathbb Z[1/2]\otimes_{\mathbb Z}\mathbb Z/2\mathbb Z=0,
	\]
	所以不忠实平坦。平坦模局部自由的定理需要有限呈示，此例不满足：若有限多个分母分别为二的幂，取其中最大指数后，它们的整数线性组合不能产生更高次分母的倒数。
\end{solution}

\begin{exercise}\label{g5:ex:fitting-triangular-family}
	设 \(k\) 为域，\(A=k[s,t]\)，令 \(M\) 由呈示
	\[
		A^2\xrightarrow{\left(\begin{smallmatrix}s&t\\0&s\end{smallmatrix}\right)}
		A^2\longrightarrow M\longrightarrow0
	\]
	定义。计算全部 Fitting 理想、纤维维数分层及 \(M\) 的平坦位置。
\end{exercise}
\begin{solution}
	子式直接给出
	\[
		\operatorname{Fitt}_0(M)=(s^2),\qquad
		\operatorname{Fitt}_1(M)=(s,t),\qquad
		\operatorname{Fitt}_i(M)=A\quad(i\geq2).
	\]
	在 \(D(s)\) 上矩阵可逆，纤维为零；在 \(V(s)\cap D(t)\) 上矩阵秩一，纤维维数一；在原点 \((s,t)\) 上矩阵为零，纤维维数二。这些集合覆盖素谱。

	\(D(s)\) 上 \(M\) 本身为零，故平坦。若 \(\mathfrak p\) 包含 \(s\) 而不含 \(t\)，局部自由的唯一可能秩为一，但 \((s^2)A_{\mathfrak p}\ne0\)，不满足该秩所需的 \(\operatorname{Fitt}_0=0\)。原点处唯一可能秩为二，但 \((s,t)A_{(s,t)}\ne0\)，不满足 \(\operatorname{Fitt}_1=0\)。因此平坦位置恰为 \(D(s)\)。这里用到了有限呈示平坦模的局部自由刻画；纤维矩阵降秩本身没有使局部环中的子式变成零。
\end{solution}

\begin{exercise}\label{g5:ex:fitting-direct-sum}
	设 \(M,N\) 为交换环 \(A\) 上的有限呈示模，证明对 \(r\geq0\)，
	\[
		\operatorname{Fitt}_r(M\oplus N)
		=\sum_{i+j=r}\operatorname{Fitt}_i(M)\operatorname{Fitt}_j(N),
		\qquad i,j\geq0.
	\]
\end{exercise}
\begin{solution}
	分别取具有 \(m,n\) 行的呈示矩阵 \(U,V\)。直和的呈示矩阵为分块对角矩阵 \(\operatorname{diag}(U,V)\)。其一个 \((m+n-r)\) 阶子式若非零，来自两个对角块的所选行数与列数必须分别相等，否则每个行列式展开项都包含一个零块中的系数。这样的子式为一个 \(a\) 阶 \(U\) 子式乘一个 \(b\) 阶 \(V\) 子式，至多相差符号，其中 \(a+b=m+n-r\)。

	令 \(i=m-a\)、\(j=n-b\)，便有 \(i+j=r\)，所得乘积属于右端。反过来，若 \(0\leq i\leq m\)、\(0\leq j\leq n\)，分别从两块选定子式，将它们放在同一分块子矩阵中，乘积就是原矩阵的一个子式。若 \(i>m\) 且 \(r<m+n\)，则 \(j=r-i<r-m\)，Fitting 理想随下标递增，故该项包含于下标 \((m,r-m)\) 对应的项中；\(j>n\) 的情形同理。这就得到反向包含。若 \(r\geq m+n\)，左端为单位理想，右端也含有两项均为单位理想的乘积；零阶子式取 \(1\)、超过矩阵可选阶数的子式取零，与上述计算一致。
\end{solution}

\begin{exercise}\label{g5:ex:nonfinite-flat-locus}
	设 \(k\) 为无限代数闭域，\(A=k[t]\)，令
	\[
		M=\bigoplus_{a\in k}A/(t-a).
	\]
	求出 \(M\) 的平坦位置，并证明这个集合不是开集。说明它为何不违背本章的开放性结论。
\end{exercise}
\begin{solution}
	在泛型点 \((0)\) 处，每个非零多项式 \(t-a\) 都变为单位，因此各直和分量的局部化都为零。局部化与直和可交换，故 \(M_{(0)}=0\)，它在分式域上平坦。

	在闭点 \(\mathfrak m_b=(t-b)\) 处，当 \(a\ne b\) 时，\(t-a\notin\mathfrak m_b\)，相应分量消失；当 \(a=b\) 时，分量变为剩余域。因此
	\[
		M_{\mathfrak m_b}\cong A_{\mathfrak m_b}/(t-b).
	\]
	这个非零模被局部整环中的非零元素 \(t-b\) 杀死，不是无挠模，因而不平坦。\(k\) 代数闭保证素谱只有泛型点与这些闭点，所以平坦位置恰为 \(\{(0)\}\)。

	若这个单点集开放，就包含泛型点的某个主开邻域 \(D(f)\)，其中 \(f\ne0\)。非零多项式只有有限多个根，而 \(k\) 无限，故可取 \(b\) 使 \(f(b)\ne0\)。于是闭点 \(\mathfrak m_b\in D(f)\)，矛盾。

	模 \(M\) 甚至不有限生成：有限多个元素的直和支撑只涉及有限多个参数，其生成子模不可能覆盖其余非零分量。因此它不满足\cref{g5:prop:finite-module-flat-locus}的有限呈示假设。这个例子也说明，在泛型点处已经自由，若缺少有限性，仍未必能把自由性扩展到任何开邻域。
\end{solution}

\begin{exercise}\label{g5:ex:generic-free-empty-fiber}
	设 \(k\) 为域，\(A=k[t]\)，\(B=A[x]\)，\(M=B/(tx-1)\)。证明 \(M\) 是有限生成 \(B\) 模并在 \(A\) 上平坦；计算它在全部素点处的纤维维数，判断它是否为有限呈示 \(A\) 模。
\end{exercise}
\begin{solution}
	\(M\) 作为 \(B\) 模由 \(1\) 生成。环同构 \(B/(tx-1)\cong A_t\) 将 \(x\) 送到 \(t^{-1}\)，故它在 \(A\) 上平坦。在不含 \(t\) 的每个素点，纤维为 \(\kappa(\mathfrak p)\)，维数一；在 \((t)\) 处，关系变为 \(-1=0\)，纤维为零。因此纤维维数在 \((t)\) 附近不恒定。

	这个 \(A\) 模甚至不有限生成：有限多个有理函数 \(a_i(t)/t^{n_i}\) 的分母有共同上界，而它们的 \(A\) 线性组合不能给出任意更高幂的 \(t^{-1}\)。故不能将有限呈示 \(A\) 模的局部自由结论用于它。泛自由性却仍适用，因为它为有限生成 \(B\) 模；在 \(D(t)\) 上它确为秩一的自由 \(A_t\) 模。
\end{solution}

\begin{exercise}\label{g5:ex:generic-zero-module}
	设 \(A=k[t]\)、\(P=A[x]\)，令
	\[
		M=P/(tx^2-x,tx-t).
	\]
	确定 \(M\) 作为 \(A\) 模的结构、纤维和一个使它自由的非零局部化元素，并解释其泛自由秩。
\end{exercise}
\begin{solution}
	在商环中，由 \(tx=t\) 乘 \(x\) 得 \(tx^2=tx=t\)，结合 \(tx^2=x\) 得 \(x=t\)，再代回得到 \(t^2=t\)。反向代入 \(x=t\)、\(t^2=t\) 也满足两个原关系，故
	\[
		M\cong A/(t(t-1)).
	\]
	在 \((t)\)、\((t-1)\) 处纤维各为一维，在其他素点处纤维为零。取 \(s=t(t-1)\ne0\)，则 \(M_s=0\)，它是秩零的自由 \(A_s\) 模。分式域上原关系生成整个多项式环，因此标准单项式集合为空；泛自由性允许这种情形，并不预设模在泛型点非零。
\end{solution}

\begin{exercise}\label{g5:ex:faithful-descent-properties}
	设 \(A\to B\) 忠实平坦，\(M\) 为任意 \(A\) 模，记 \(M_B=B\otimes_A M\)。
	\begin{enumerate}
		\item 证明 \(M\) 平坦当且仅当 \(M_B\) 为平坦 \(B\) 模。
		\item 证明有限生成、有限呈示和有限生成投射这三个性质均可由 \(M_B\) 传回 \(M\)。
	\end{enumerate}
\end{exercise}
\begin{solution}
	1. 正方向由平坦性的基变换性质。反过来，对任意 \(A\) 模单射 \(L\to N\)，先张量 \(B\)，再在 \(B\) 上张量平坦模 \(M_B\)，得到单射
	\[
		(B\otimes_A L)\otimes_B M_B
		\longrightarrow(B\otimes_A N)\otimes_B M_B.
	\]
	它就是 \(L\otimes_A M\to N\otimes_A M\) 再张量 \(B\) 后的映射。忠实平坦性反映单射，故 \(M\) 平坦。

	2. 若 \(M_B\) 有限生成，将有限组生成元各自写成有限个纯张量之和，收集其中出现的有限多个 \(M\) 元素，使它们生成子模 \(M_0\subseteq M\)。则 \(B\otimes_A M_0\to M_B\) 满射，故 \(B\otimes_A(M/M_0)=0\)。忠实性给出 \(M=M_0\)，从而有限生成。

	若 \(M_B\) 有限呈示，先由上一步选有限满射 \(A^n\to M\)，记核为 \(K\)。平坦性给出正合列
	\(0\to B\otimes_AK\to B^n\to M_B\to0\)。
	由\cref{g5:prop:finite-presentation-kernel}，\(B\otimes_A K\) 有限生成。再用上一步得 \(K\) 有限生成，故 \(M\) 有限呈示。

	若 \(M_B\) 有限生成投射，则它有限呈示且平坦。前两步将这两个性质传回 \(M\)，再用\cref{g5:thm:finite-presented-flat}，得到 \(M\) 有限生成投射。这里有限生成投射模有限呈示，也可从它为有限自由模直和分量的表示直接看出。
\end{solution}

\begin{exercise}\label{g5:ex:integers-principal-cover}
	证明 \(\mathbb Z\to\mathbb Z[1/2]\times\mathbb Z[1/3]\) 忠实平坦，并将双重基变换的相容条件化为下面的正合列：
	\[
		0\longrightarrow\mathbb Z
		\longrightarrow\mathbb Z[1/2]\oplus\mathbb Z[1/3]
		\xrightarrow{(a,b)\mapsto a-b}\mathbb Z[1/6]
		\longrightarrow0.
	\]
\end{exercise}
\begin{solution}
	没有素理想同时包含 \(2\) 和 \(3\)，所以 \(D(2),D(3)\) 覆盖整数素谱，正文的有限主开覆盖例给出忠实平坦性。双重张量的两个交叉因子均为 \(\mathbb Z[1/6]\)，同一开集与自身的因子不增加条件，因此相容性就是 \(a=b\) 于 \(\mathbb Z[1/6]\)。

	若 \(a=b\)，这个有理数的最简分母既只能含素因子 \(2\)，又只能含素因子 \(3\)，所以为整数，证明中间正合。为证最后满射，只需处理 \(1/(2^m3^n)\)。由于 \(2^m,3^n\) 互素，可选整数 \(u,v\) 使 \(u3^n+v2^m=1\)，于是
	\[
		\frac1{2^m3^n}=\frac{u}{2^m}+\frac{v}{3^n}.
	\]
	将第二项写成差中的负项即可。这一步给出本例特有的末端满射；一般等化子序列并没有宣称其最后一张差映射满射。
\end{solution}

\begin{exercise}\label{g5:ex:cusp-node-flat-family}
	设 \(k\) 为代数闭的特征零域，取
	\[
		A=k[t],\qquad B=A[x,y]/(y^2-x^3-tx^2).
	\]
	证明 \(B\) 为忠实平坦 \(A\) 代数，计算原点处各纤维的切锥，并将 \(t=a\ne0\) 的纤维化为第四章的结点方程。
\end{exercise}
\begin{solution}
	关于 \(y\) 的方程首一，除法给出唯一余式 \(f(t,x)+y g(t,x)\)。故 \(B\) 为自由 \(A[x]\) 模，基为 \(1,y\)，也为自由 \(A\) 模，基为
	\[
		x^n,\quad x^ny\qquad(n\geq0).
	\]
	这是一组非空基，所以忠实平坦。\(t=0\) 的纤维为尖点 \(y^2=x^3\)，其切锥为 \(Y^2=0\)。当 \(t=a\ne0\) 时，最低次方程为 \(Y^2-aX^2=0\)，在 \(k\) 中分成两个不同的线性因子。

	选 \(b\in k\) 使 \(b^2=a\)，令 \(x=aX\)、\(y=abY\)。因 \(a,b\ne0\)，这是可逆线性换元，方程变为
	\[
		Y^2=X^2(X+1).
	\]
	因此非零参数给出本书已研究的结点，零参数给出尖点。此族的纤维始终为一维，原点切空间始终为二维，改变的是切锥的乘法结构和分支；这些变化与平坦性相容。
\end{solution}

\begin{exercise}\label{g5:ex:complex-faithful-equalizer}
	构造 \(\mathbb C\otimes_{\mathbb R}\mathbb C\cong\mathbb C\times\mathbb C\)，并在这一表示中验证 \(\mathbb R\) 是两张映射
	\(\mathbb C\rightrightarrows\mathbb C\otimes_{\mathbb R}\mathbb C\)
	的等化子。
\end{exercise}
\begin{solution}
	定义 \(\Phi(z\otimes w)=(zw,z\overline w)\)。它对实数平衡且保持乘法，故为环同态。把左端按第一个因子看作复向量空间，基为 \(1\otimes1,1\otimes i\)，其像为 \((1,1),(i,-i)\)，行列式为 \(-2i\ne0\)，所以 \(\Phi\) 为同构。两张映射在右端分别为
	\[
		z\otimes1\longmapsto(z,z),\qquad
		1\otimes z\longmapsto(z,\overline z).
	\]
	它们相等当且仅当 \(z=\overline z\)，也就是 \(z\in\mathbb R\)。这与 \(\mathbb R\to\mathbb C\) 为忠实平坦扩张一致：作为实模，\(\mathbb C\) 自由且基非空。谱映射两边都只有一个点，因此满射；这个事实并不要求存在实代数同态 \(\mathbb C\to\mathbb R\)，后者会要求 \(i^2=-1\) 在实数中有解，因而不存在。
\end{solution}
